\documentclass[11pt]{article}
\usepackage[T1]{fontenc}
\usepackage{lmodern,microtype,amsmath,amssymb,amsthm,mathtools,mathrsfs}
\usepackage[margin=1in]{geometry}
\usepackage{booktabs,longtable,array,enumitem}
\usepackage{xcolor}
\usepackage[colorlinks=true,linkcolor=blue!45!black,citecolor=blue!45!black,urlcolor=blue!45!black]{hyperref}
\usepackage[nameinlink,noabbrev]{cleveref}
\hypersetup{pdftitle={Finite Fluctuation Hierarchies for the Negative-Time Riemann Heat Flow},
pdfauthor={Edward Baker},
pdfsubject={Canonical approximation, effective evaluation, and explicit bounds}}
\setlength{\emergencystretch}{2em}
\allowdisplaybreaks[2]
\numberwithin{equation}{section}
\newtheorem{theorem}{Theorem}[section]
\newtheorem{proposition}[theorem]{Proposition}
\newtheorem{lemma}[theorem]{Lemma}
\newtheorem{corollary}[theorem]{Corollary}
\theoremstyle{definition}
\newtheorem{definition}[theorem]{Definition}
\newtheorem{assumption}[theorem]{Classical input}
\theoremstyle{remark}
\newtheorem{remark}[theorem]{Remark}
\newcommand{\dd}{\,\mathrm d}
\newcommand{\ii}{\mathrm i}
\newcommand{\e}{\mathrm e}
\newcommand{\R}{\mathbb R}
\newcommand{\C}{\mathbb C}
\newcommand{\N}{\mathbb N}
\newcommand{\Rea}{\operatorname{Re}}
\newcommand{\Ima}{\operatorname{Im}}
\DeclareMathOperator{\Log}{Log}
\newcommand{\Lp}{\log^{+}}
\newcommand{\calQ}{\mathcal Q}
\newcommand{\calH}{\mathscr H}
\newcommand{\calE}{\mathcal E}
\newcommand{\normg}[1]{\lVert #1\rVert_{\gamma;X;\tau,B}^{(3,1)}}
\newcommand{\normo}[1]{\lVert #1\rVert_{\mathrm o;X;1,2}^{(3,1)}}
\newcommand{\normc}[1]{\lVert #1\rVert_{\mathrm c;X;1,2}^{(3,1)}}
\newcommand{\poly}[2]{#1[#2]}
\newcommand{\abs}[1]{\lvert #1\rvert}
\newcommand{\bigabs}[1]{\left|#1\right|}
\title{Finite Fluctuation Hierarchies for the\\Negative-Time Riemann Heat Flow\\[0.5em]
\large Canonical approximation, effective evaluation,\\
and explicit bounds}
\author{Edward Baker\thanks{Draft prepared for Edward Baker with substantial assistance from OpenAI language models. The underlying analytical and computational notes were developed with gpt-6.1-sol (reasoning effort ultra). This synthesis, including the finite algorithms, calibration and cost results, was prepared with gpt-6-astra (recorded effort max). The additive-cost onset and uniform RS-bound strengthenings were developed and integrated with gpt-6.1-sol (actual effort ultra). Internal LLM assessments are not independent specialist validation.}}
\date{26 October 2026\\\small Working manuscript}
\begin{document}
\maketitle
\begin{abstract}
We construct a finite fluctuation hierarchy for the canonical negative-time
Riemann heat flow \(h(T,z)=H(-T,z)\), \(T>0\). Gaussian moments evolve
the full retained arithmetic population, while a bounded integral preserves
its real Riemann--Siegel completion. Grading by logarithmic curvature gives
an exact finite recurrence. At fixed order \(p\), completion through
\(C_{2p}\) yields natural gamma-scale error \(O(X^{-(p+3/4)})\),
simultaneously for the field, three spatial derivatives, and one time
derivative on fixed complex charts and compact positive time intervals.
Controlling the order-dependent constants before choosing
\(p=\lfloor\tfrac12\log(X/(4\pi))\rfloor\) gives outer-envelope error
\(\exp[-\tfrac14\log^2(X/(4\pi))]\) on \(1\le T\le2\).
We supply finite algorithms for every required coefficient, scale,
quadrature datum, and frozen partition, and prove explicit uniform
calibration constants and an all-height sufficient log-height
\(L\ge64(\lceil r+b+2\rceil+1)\), retaining the earlier quartic bound
as a fallback. A separate moving-endpoint argument gives the same total
accuracy on \(L^{-1}\le T\le2\) at sufficient
\(L\ge1024(\lceil r+b+2\rceil+1)^2\), with complete endpoint,
population and tail costs. On certified
dyadic-center, strictly witnessed rational-query inputs, these constructions
preserve the same total field and derivative accuracy. The exact moving
endpoint also has a literal-tag interface with a paid time-input error.
Its sufficient outer quadrature order is \(O(L^3)\), separately from
provider-internal work. Sufficient delivered precision
is \(3\log^2(X/(4\pi))/(8\log2)+O(\log X\log\log X)\) bits;
bit work is the full arithmetic population times a polynomial in input
length, precision, log-height, and center/time admissibility separation, with polynomial
streamed storage. The onset and cost bounds are conservative.
A uniform boundary argument reduces every coefficient majorant by
a factor greater than ten, and a separate contour bound permits eleven
fewer Gauss nodes per edge at identical jet accuracy.
Finite component replay is supporting evidence; a complete certified
implementation, strong negative-time precision through zero, and transfer
to specified physical observables remain open.
\end{abstract}

\clearpage
\begingroup
\small
\tableofcontents
\endgroup
\clearpage
\section{Introduction and the approximation choices}
\label{sec:intro}
The useful question for a finite asymptotic formula is what its
retained terms buy. Here there are two separate resolutions: the
curvature of the arithmetic carriers and the order of their
Riemann--Siegel completion. Improving either in isolation need not
improve the approximation to the canonical field. This paper derives
their matching, then controls the constants sufficiently well to let
the chosen order increase with height.

The development synthesizes the comparison and evaluation results in
recurrence Notes 02--06, 09--11, 16 and 18
\cite{rec02,rec03,rec04,rec05,rec06,rec09,rec10,rec11,rec16,rec18}. The comparison is always to the
same complete theta-source heat flow. In particular, a finite
polynomial hierarchy is not a truncation of the theta source or
of the retained arithmetic population. Its exact heat evolution
is useful structure, but the error is proved by matching initial
data, local remainders, and tails.

\subsection{What survives from the preceding ansatz investigation}
\label{subsec:retention}
The preceding investigation separated the elementary outer carrier
from the complete short arithmetic approximation
\cite{ans03,ans04,ans05,ans10,ans12}. We retain that distinction.
Put \(L(z)=\Log(z/(4\pi))\), with the branch continued from the
positive real axis, and define
\begin{subequations}\label{eq:outer-arms}
\begin{align}
 A_{\mathrm o}(t,z)&=\frac{\pi^{1/4}}{16}z^{7/4}
 \exp\left\{-\frac{\pi z}{8}
       +\frac t{16}\left(L(z)^2-\frac{\pi^2}{4}\right)\right\},\\
 \Theta_{\mathrm o}(t,z)&=\frac z4(L(z)-1)
                  +\frac{\pi t}{16}L(z)+\frac{7\pi}{8}.
\end{align}
\end{subequations}
The elementary cosine \(A_{\mathrm o}\cos\Theta_{\mathrm o}\)
describes the first carrier. Its adequacy is regime-dependent.
In the negative-time regime studied here, the full arithmetic
population cannot be replaced by that carrier using the present
estimates.

For a fixed chart cutoff \(N=\lfloor\sqrt{X/(4\pi)}\rfloor\), the
earlier simplest complete short formula is
\begin{subequations}\label{eq:short}
\begin{align}
 F_{\mathrm{short}}(t,z)
 &=A_{\mathrm o}(t,z)\sum_{n=1}^{N}w_n(t,z)\cos\phi_n(t,z)
       +\e^{-\pi^2t/64}\mathcal R_N(z),\\
 w_n(t,z)&=n^{-1/2}
     \exp\left\{\frac t4\bigl((\log n)^2-L(z)\log n\bigr)\right\},\\
 \phi_n(t,z)&=\Theta_{\mathrm o}(t,z)
                  -\left(\frac z2+\frac{\pi t}{8}\right)\log n,\\
 \mathcal R_N(z)&=(-1)^NG(z)a(z)^{-1/2}\mathcal C(a(z)-N),\\
 \mathcal C(v)&=\frac{\cos(2\pi(v^2-v)-\pi/8)}{\cos(2\pi v)},
 \qquad a(z)=\sqrt{z/(4\pi)}.
\end{align}
\end{subequations}
Here \(G\) is defined in \cref{eq:gamma-scale}, and removable
singularities in \(\mathcal C\) are filled analytically.
The cutoff is frozen when differentiating. Ansatz Note 05 proves
natural-scale \(O(X^{-3/4})\) field, spatial-through-three, and
first-time errors on \(1\le T=-t\le2\).
Ansatz Note 04 separately proves its signed crossover estimates on
\(\abs t\log(X/(4\pi))\le8\). These are different domains and
different weighted statements.

The first higher approximation replaces the elementary arithmetic
carriers by exact gamma modes with a curvature correction. Its
Riemann--Siegel completion must also improve: the earlier
target-dependent construction retains \(C_0,C_1\) in a bounded
real-data heat integral and has \(O(X^{-5/4})\) error. The present
frozen first level retains \(C_0,C_1,C_2\) and has
\(O(X^{-7/4})\) error. The source of this difference will be
derived in \cref{subsec:matching,subsec:low-orders}.

\begin{center}
\small
\begin{tabular}{@{}>{\raggedright\arraybackslash}p{.21\textwidth}>{\raggedright\arraybackslash}p{.29\textwidth}>{\raggedright\arraybackslash}p{.24\textwidth}>{\raggedright\arraybackslash}p{.14\textwidth}@{}}
\toprule
Approximation & Arithmetic terms retained & Completion retained &
Natural-scale error\\
\midrule
Earlier short formula &
All \(n\le N\), both elementary arms &
Leading \(\mathcal R_N\), with its heat factor &
\(X^{-3/4}\)\\[3pt]
Earlier first curvature correction &
All exact gamma modes, target-dependent quadratic correction &
Bounded integral through \(C_1\), exact transfer &
\(X^{-5/4}\)\\[3pt]
Frozen first level &
\(U_0+U_1\), all modes &
Bounded integral through \(C_2\) &
\(X^{-7/4}\)\\[3pt]
Frozen second level &
\(U_0+U_1+U_2\), all modes &
Bounded integral through \(C_4\) &
\(X^{-11/4}\)\\[3pt]
Frozen order \(p\ge1\) &
\(\sum_{j=0}^pU_j\), all modes &
Bounded integral through \(C_{2p}\) &
\(X^{-(p+3/4)}\)\\
\bottomrule
\end{tabular}
\end{center}
All rows apply on \(1\le T\le2\) with fixed buffers and sufficiently
large height; the frozen fixed-order rows extend to any declared
fixed compact interval \(0<\tau\le T\le B\).
Every integral row includes the exact change between the natural
cutoff at the integration point and the chart cutoff \(N\).
The table compares proved upper allowances, not pointwise
monotonicity of approximation error.

\subsection{Main conclusions and proof architecture}
The paper develops three connected conclusions. First,
\cref{thm:fixed,thm:growing} quantify the gain from each retained
fluctuation level, first at fixed order and then at controlled growing
order. The completion order must track the arithmetic curvature order:
improving one without the other can leave the error rate unchanged.

Second, \cref{thm:effective} turns this comparison into a terminating
finite evaluation construction. Analytic coefficient bounds, quadrature,
input sensitivity, and roundoff are allocated separately. Shifted gamma
expansions, exact Gauss root counts, pole-safe coefficient formulas, and
interior branch guards then supply every data contract. The accepted
input class consists of dyadic centers with strict certified descriptors
and exact rational queries; its separation dependence is explicit.

Third, \cref{thm:linear-onset,thm:bit-cost} make the construction
effective in height and cost. With \(\nu=\lceil r+b+2\rceil+1\),
a sufficient log-height is \(L_{\rm new}=64\nu\), improving the earlier
\(L_*=2^{22}\nu^4\) by the factor \(2^{16}\nu^3\). The bit-work bound
retains all \(N\asymp\sqrt X\) modes, while streaming avoids storing
them all. These conservative bounds establish a finite procedure; their
size does not assert practical evaluation at the sufficient onset.

The proofs proceed from the real initial-value heat formula to the local
generating function. Integrating its weighted remainder against Gaussian
moments preserves the useful powers of height; the full population costs
\(X^{1/4}\), and matching the initial Riemann--Siegel remainder fixes
\(K=2p\). True and artificial tails are kept separate. Only after their
order dependence is controlled do we change to the smaller outer envelope.
The evaluation budget is then proved as a modular theorem, supplied by
finite algorithms, and closed with explicit uniform calibration.

The analytical claims are proved from the named classical inputs.
Their source notes have same-model internal assessments, and this
manuscript provides an internal synthesis check. Supporting component
computations are linked in \cref{sec:effective}; they are not premises of
the theorems. Independent specialist validation, including the source
proofs behind the imported coefficient identities, remains open.
No claim of priority, infinite-series convergence, or a conclusion about
zeros from approximation accuracy alone is made.

\section{Canonical heat evolution and the real initial-data completion}
\label{sec:source}
\subsection{Normalization and exact evolution}
Use the classical normalization \cite{polymath}
\begin{subequations}\label{eq:canonical}
\begin{align}
 \Phi(v)&=\sum_{n=1}^{\infty}
 (2\pi^2n^4\e^{9v}-3\pi n^2\e^{5v})\e^{-\pi n^2\e^{4v}},\\
 H(t,z)&=\int_0^\infty \e^{tv^2}\Phi(v)\cos(zv)\dd v,
 & H(0,z)&=\frac18\xi\left(\frac{1+\ii z}{2}\right).
\end{align}
\end{subequations}
The source is integrable and decays faster than a Gaussian.
Differentiation gives \(H_t=-H_{zz}\).
We write \(h(T,z)=H(-T,z)\), so \(h_T=h_{zz}\).
Let
\[
 k_T(v)=\frac1{\sqrt{4\pi T}}\exp\left(-\frac{v^2}{4T}\right),
 \qquad T>0.
\]
\begin{lemma}[Real initial-value formula]\label{lem:heat}
For every \(T>0\) and \(z\in\C\),
\begin{equation}
 h(T,z)=\int_\R k_T(z-u)H(0,u)\dd u.
 \label{eq:heat-convolution}
\end{equation}
The integral and its spatial and time derivatives converge uniformly
on compact subsets with \(T>0\).
\end{lemma}
\begin{proof}
Insert the cosine transform of \(H(0,u)\).
For real \(u\), the cosine is bounded, while
\[
 \abs{k_T(z-u)}
 =\frac{\e^{(\Ima z)^2/(4T)}}{\sqrt{4\pi T}}
      \e^{-(u-\Rea z)^2/(4T)}.
\]
Integrability of \(\Phi\) permits Fubini. The Fourier transform of
the Gaussian is \(\e^{-Tv^2+\ii zv}\), giving
\eqref{eq:heat-convolution}. Derivatives introduce fixed polynomial
factors and inverse powers of \(T\), still integrable on the stated
compact sets.
\end{proof}
This identity is the reason real Riemann--Siegel data can be used
without asserting a complex remainder theorem: the heat kernel
produces the holomorphic field.

\subsection{Charts, gamma scale, and frozen data}
Fix \(r,b>0\), \(0<\tau\le B<\infty\), and
\[
 \Omega_X(r,b)=\{z:\abs{\Rea z-X}\le r,\ \abs{\Ima z}\le b\},
 \quad R=r+1,\quad Y=b+1,\quad W=R+Y.
\]
The enlarged chart is \(\Omega_X(R,Y)\).
For physical applications the inherited choice is
\(r=\ell+2\), \(b=1+h_{\mathrm{av}}+2\); no physical estimate is
deduced here.
Throughout put
\[
 L=\log(X/(4\pi)),\quad a_X=\sqrt{X/(4\pi)},\quad
 N=\lfloor a_X\rfloor,\quad
 \bar B=\max(1,B),\quad V_B=8\sqrt{\bar B}\,L.
\]
The center \(X\), the orders, all coefficients, all real partitions,
and the endpoints \(X\pm V_B\) are fixed under both \(z\)- and
\(T\)-differentiation.

Define
\begin{subequations}\label{eq:gamma-scale}
\begin{align}
 w_\epsilon(z)&=(1+\epsilon\ii z)/2,\qquad \epsilon=\pm1,\\
 M(w)&=\frac{w(w-1)}{16}\pi^{-w/2}\Gamma(w/2),\\
 G(z)&=\{M(w_+(z))M(w_-(z))\}^{1/2},\qquad g=G(X)>0.
\end{align}
\end{subequations}
The square root is positive on the real segment and holomorphic
on sufficiently high fixed charts. Sectorial gamma estimates
\cite{dlmfgamma} give
\begin{equation}
 \frac{G'(z)}{G(z)}=-\frac{\pi}{8}+O(X^{-1}),\qquad
 G(X)\ge g_-X^{7/4}\e^{-\pi X/8},\qquad
 C_\chi\ge\sup_{\Omega_X(r,b)}\frac{g}{\abs{G(z)}},
 \label{eq:gamma-comparison}
\end{equation}
with fixed positive constants after a fixed onset.
For the growing-order slab, \cref{sec:calibration} supplies their explicit
values, full complex domains, and a common sufficient center onset.
Here and below classical constants may depend on fixed sectors
and chart radii, but not on an order subsequently allowed to grow.

For a field \(E\), the natural comparison norm is
\begin{equation}
 \normg{E}=\sup_{\substack{\tau\le T\le B\\z\in\Omega_X(r,b)}}
 \frac{\max_{0\le j\le3}\abs{\partial_z^jE(T,z)}
                      +\abs{\partial_TE(T,z)}}{\abs{G(z)}}.
 \label{eq:natural-norm}
\end{equation}
The center norm replaces \(\abs{G(z)}\) by \(g\); later the outer
norm replaces it by a positive amplitude envelope.
None divides by the field itself.

\subsection{Classical real Riemann--Siegel input}
\label{subsec:rs-input}
For real \(u\) in \(I_{X,B}=[X-V_B,X+V_B]\), write
\[
 a(u)=\sqrt{u/(4\pi)},\quad m(u)=\lfloor a(u)\rfloor,\quad
 \chi(u)=1-2(a(u)-m(u)),
\]
\[
 U(u)=\exp\left[-\ii\left\{\frac u4\log\frac u{4\pi}
                                   -\frac u4-\frac{\pi}{8}\right\}\right].
\]
The fractional coordinate \(\chi\) is distinct from the fluctuation
order \(p\). Set
\begin{subequations}\label{eq:real-completion}
\begin{align}
 S_m(u)&=\sum_{\epsilon=\pm1}\sum_{n=1}^m
                 M(w_\epsilon(u))n^{-w_\epsilon(u)},\\
 R_{m,K}(u)&=2\Rea\left\{(-1)^{m-1}M(w_+(u))U(u)a(u)^{-1/2}
                 \sum_{k=0}^{K}C_k(\chi(u),1/2)a(u)^{-k}\right\},\\
 f_K(u)&=S_{m(u)}(u)+R_{m(u),K}(u).
\end{align}
\end{subequations}
The coefficient normalization is the complex one in Polymath,
Proposition 6.2, equations (53)--(56) \cite{polymath}.
At \(\sigma=1/2\), for every positive integer \(K\), its
real initial-data remainder gives
\begin{equation}
 \abs{H(0,u)-f_K(u)}
 \le\frac{2^{7/4}}7G(u)a(u)^{-1/2}
       \Gamma((K+1)/2)\left(\frac{1.1}{a(u)}\right)^{K+1}.
 \label{eq:rs-remainder}
\end{equation}
Only this real initial formula is imported. No positive-time
evolved remainder is continued to negative time.

The complete change of cutoff is retained by setting
\begin{subequations}\label{eq:retained-integral}
\begin{align}
 Q_{X,K}(u)&=S_{m(u)}(u)-S_N(u)+R_{m(u),K}(u),\\
 \mathcal I_{X,K,B}(T,z)&=\int_{I_{X,B}}k_T(z-u)Q_{X,K}(u)\dd u.
\end{align}
\end{subequations}
If \(a(X+V_B)-a(X-V_B)<1\), there is at most one threshold
\(4\pi n^2\) in the window, and at most two real branches.
Their endpoints are independent of \(z,T\). Possible values at a
threshold have measure zero. The real-part operation remains inside
the initial data in \eqref{eq:real-completion}; moving it outside
the complex-kernel integral would change the function.

For completeness, the classical real coefficient estimate
\(\abs{C_k}\le (3\sqrt2/(2\pi))\Gamma(k/2)/2^k\), \(k\ge1\),
and \(\abs{C_0}\le1/2\), imply a bounded retained coefficient sum
when \(K\le a_{\min}^2\), where
\(a_{\min}^2=(X-V_B)/(4\pi)\). Indeed successive terms of each
parity have ratio \(k/(8a(u)^2)\le1/8\).
Together with the at-most-one exact transfer this gives
\(\abs{Q_{X,K}(u)}\le C G(u)X^{-1/4}\).
This is a size bound, not a regularity or quadrature theorem.

\subsection{The smaller outer envelope}
Let \(C_\pm(t,z)=A_{\mathrm o}(t,z)\e^{\pm\ii\Theta_{\mathrm o}(t,z)}/2\)
and \(\calE_{\mathrm o}(t,z)=\abs{C_+}+\abs{C_-}\).
For \(1\le T\le2\), gamma calibration yields
\begin{equation}
 \calE_{\mathrm o}(-T,z)\ge
 c_{\mathrm o}\abs{G(z)}
    \exp\left(-\frac{TL^2}{16}+\frac{\pi^2T}{64}\right).
 \label{eq:outer-calibration}
\end{equation}
One way to see the scale is to compare the modulus of
\(A_{\mathrm o}(0,z)\) with \(\abs{G(z)}\) by sectorial Stirling.
On a fixed chart \(L(z)=L+O(X^{-1})\), so the heat factor changes
by \(\exp(O((1+L)/X))\). Finally
\((\abs{\e^{\ii\Theta}}+\abs{\e^{-\ii\Theta}})/2\ge1\).
After increasing the onset, fixed comparison constants give
\eqref{eq:outer-calibration}, as in \cite{ans03,ans05}.
Thus
\begin{equation}
 \normo{E}\le c_{\mathrm o}^{-1}\e^{L^2/8}
       \lVert E\rVert_{\gamma;X;1,2}^{(3,1)}.
 \label{eq:outer-conversion}
\end{equation}
The factor \(\e^{L^2/8}\) is why fixed algebraic error estimates
alone do not establish the desired outer accuracy.

\section{The finite fluctuation hierarchy}
\label{sec:hierarchy}
\subsection{A common logarithmic shape}
On the high-height disk choose an analytic logarithm
\(\ell_\epsilon(z)=\log M(w_\epsilon(z))\), and put
\begin{subequations}\label{eq:frozen-data}
\begin{align}
 b_{\epsilon n}&=M(w_\epsilon(X))n^{-w_\epsilon(X)},&
 \lambda_{\epsilon n}&=\ell_\epsilon'(X)-\epsilon\ii\log n/2,\\
 c_{\epsilon r}&=\frac{\ell_\epsilon^{(r+1)}(X)}{(r+1)!},&
 A_\epsilon(q)&=\exp\{\ell_\epsilon(X+q)-\ell_\epsilon(X)
                                      -\ell_\epsilon'(X)q\}.
\end{align}
\end{subequations}
Then, exactly for local real \(q\),
\[
 M(w_\epsilon(X+q))n^{-w_\epsilon(X+q)}
       =b_{\epsilon n}\e^{\lambda_{\epsilon n}q}A_\epsilon(q).
\]
All modes in one arm share the same shape \(A_\epsilon\).
This separates arithmetic frequency from gamma curvature.

\begin{lemma}[Simultaneous coefficient bound]\label{lem:coeff}
There are fixed \(A,H\ge1\), \(D\ge4\), and a fixed onset such that
\begin{equation}
 \sup_{\abs{q}\le X/4}\abs{\ell_\epsilon''(X+q)}\le H/X,
 \qquad
 \abs{c_{\epsilon r}}\le
 \frac{H4^{r-1}}{r(r+1)}X^{-r}\le AD^{r-1}X^{-r}
 \quad(r\ge1).
 \label{eq:coefficient-bound}
\end{equation}
\end{lemma}
\begin{proof}
The arguments \(w_\epsilon(z)/2\) stay in fixed sectors away from
the negative real ray and have modulus comparable to \(X\).
Sectorial digamma asymptotics and an interior Cauchy estimate give
\(\psi'(v)=1/v+O(\abs v^{-2})\). Direct differentiation gives
\[
 \ell_\epsilon''(z)=\frac1{4w_\epsilon(z)^2}
 +\frac1{4(w_\epsilon(z)-1)^2}
 -\frac1{16}\psi'(w_\epsilon(z)/2)=O(X^{-1}).
\]
Cauchy's estimate for the \((r-1)\)-st derivative of
\(\ell_\epsilon''\), divided by \((r+1)!\), proves
\eqref{eq:coefficient-bound}. The sector can first be slightly
enlarged to justify the estimate on the closed disk.
\end{proof}
This analytic disk bound is stronger than a list of fixed-\(r\)
asymptotic statements: it controls all coefficient orders at once.

\subsection{Grading by fluctuation weight}
Suppress \(\epsilon\) temporarily. Define the finite polynomials
\begin{subequations}\label{eq:partition}
\begin{align}
 a_0(q)&=1,\\
 a_j(q)&=\sum_{\substack{k_1,\ldots,k_j\ge0\\\sum_{r=1}^j rk_r=j}}
           \prod_{r=1}^j\frac{(c_rq^{r+1})^{k_r}}{k_r!},
 &P_p(q)&=\sum_{j=0}^p a_j(q).
\end{align}
\end{subequations}
For \(j\ge1\), the degrees of \(a_j\) run from \(j+1\) to \(2j\).
The index \(j\) is the power of \(X^{-1}\) suggested by
\eqref{eq:coefficient-bound}, not ordinary polynomial degree.
Locally the convergent generating function is
\[
 \exp\left(\sum_{r\ge1}c_rq^{r+1}s^r\right)=\sum_{j\ge0}a_j(q)s^j.
\]
Differentiating with respect to \(s\) gives the exact recurrence
\begin{equation}
 j a_j(q)=\sum_{r=1}^j r c_rq^{r+1}a_{j-r}(q).
 \label{eq:initial-recurrence}
\end{equation}
The first levels are
\begin{equation}
 P_1=1+c_1q^2,\qquad
 P_2=1+c_1q^2+c_2q^3+\tfrac12c_1^2q^4.
 \label{eq:first-polynomials}
\end{equation}
A degree-four Taylor polynomial would also include \(c_3q^4\),
which has fluctuation weight three. The two truncations are
therefore different even at low order.

\subsection{Gaussian moments and the evolved recurrence}
Let \(d=z-X\) and \(\calQ_T=d+2T\partial_d\).
The kernel identity \(\calQ_Tk_T(d-q)=qk_T(d-q)\) gives
\begin{subequations}\label{eq:moment}
\begin{align}
 \int_\R k_T(d-q)\e^{\lambda q}q^m\dd q
 &=\e^{\lambda d+T\lambda^2}\calH_m(d+2T\lambda,T),\\
 \calH_m(t,T)&=m!\sum_{\nu=0}^{\lfloor m/2\rfloor}
            \frac{T^\nu t^{m-2\nu}}{\nu!(m-2\nu)!}.
\end{align}
\end{subequations}
Equivalently,
\(\sum_m\calH_m(t,T)s^m/m!=\e^{ts+Ts^2}\), whence
\begin{equation}
 \calH_0=1,\quad\calH_1=t,\quad
 \calH_m=t\calH_{m-1}+2T(m-1)\calH_{m-2}.
 \label{eq:moment-recurrence}
\end{equation}
Define
\[
 U_{\epsilon,j}(T,d)
  =\sum_{n=1}^N\int_\R k_T(d-q)b_{\epsilon n}
                         \e^{\lambda_{\epsilon n}q}a_{\epsilon,j}(q)\dd q.
\]
All these integrals are absolutely convergent. Evolving
\eqref{eq:initial-recurrence} yields
\begin{equation}
 U_{\epsilon,0}=\sum_{n=1}^Nb_{\epsilon n}
                      \e^{\lambda_{\epsilon n}d+T\lambda_{\epsilon n}^2},
 \qquad
 jU_{\epsilon,j}=\sum_{r=1}^j
                  r c_{\epsilon r}\calQ_T^{r+1}U_{\epsilon,j-r}.
 \label{eq:evolved-recurrence}
\end{equation}
There is no noncommutative ordering ambiguity: all operators on the
right are powers of the same \(\calQ_T\).

\begin{definition}[Finite completed model]\label{def:model}
For finite \(p,K\), let
\begin{equation}
 F_{X,B}^{[p,K]}(T,z)=
       \sum_{\epsilon=\pm1}\sum_{j=0}^p U_{\epsilon,j}(T,z-X)
                    +\mathcal I_{X,K,B}(T,z).
 \label{eq:model}
\end{equation}
\end{definition}
\begin{proposition}[Exact frozen evolution]\label{prop:residual}
Every model \eqref{eq:model} satisfies
\[
 (\partial_T-\partial_z^2)F_{X,B}^{[p,K]}=0,\qquad
 \partial_z(\partial_T-\partial_z^2)F_{X,B}^{[p,K]}=0.
\]
\end{proposition}
\begin{proof}
Either differentiate the Gaussian transforms or use
\([\,\partial_T-\partial_d^2,\calQ_T\,]=0\).
The base exponentials solve heat exactly. The retained integral
also does, because its data, partitions, and endpoints are fixed.
\end{proof}
These identities concern the model's evolution. They do not yet
compare its initial data with the canonical source.

\subsection{A local remainder with explicit order dependence}
For a polynomial \(P(x)=\sum_m P_mx^m\) and a sequence \(M_m\), write
\(P[M]=\sum_mP_mM_m\). Introduce the positive polynomials
\begin{subequations}\label{eq:majorant-polynomials}
\begin{align}
 \mathsf C_j(x)&=\sum_{k=1}^j
     \frac{A^kD^{j-k}}{k!}\binom{j-1}{k-1}x^{j+k},
 &\mathsf H_p(x)&=1+\sum_{j=1}^p\mathsf C_j(x),\\
 \mathsf B_m(x)&=\sum_{k=1}^m
       \frac{2^{m+k}A^kD^{m-k}}{k!}\binom{m-1}{k-1}x^{m+k}.
\end{align}
\end{subequations}
\begin{lemma}[Weighted remainder]\label{lem:weighted}
If \(DV_B/X\le1/2\) and \(2AV_B^2/X\le1\), then for every integer
\(p\ge0\) and \(\abs q\le V_B\),
\begin{equation}
 \abs{A_\epsilon(q)-P_{\epsilon,p}(q)}
      \le \e X^{-(p+1)}\mathsf B_{p+1}(\abs q).
 \label{eq:weighted-remainder}
\end{equation}
Also \(\abs{a_{\epsilon,j}(q)}\le X^{-j}\mathsf C_j(\abs q)\),
and, for \(X\ge1\), \(\abs{P_{\epsilon,p}(q)}\le\mathsf H_p(\abs q)\)
on the whole real line.
\end{lemma}
\begin{proof}
Grouping a partition by \(k=\sum k_r\) uses
\[
 \sum_{\substack{\sum rk_r=j\\\sum k_r=k}}\frac1{\prod k_r!}
                 =\frac1{k!}\binom{j-1}{k-1}.
\]
This follows by taking the coefficient of \(t^j\) in
\((t+t^2+\cdots)^k/k!\), and proves the coefficient majorant.
For the remainder set
\(f(s)=\sum_{r\ge1}c_rq^{r+1}s^r\), \(x=\abs q\), \(0\le s\le1\).
The geometric majorant gives \(\abs{f(s)}\le1\) and
\[
 \frac{\abs{f^{(r)}(s)}}{r!}
 \le AD^{r-1}X^{-r}x^{r+1}(1-Dsx/X)^{-r-1}
 \le2^{r+1}AD^{r-1}X^{-r}x^{r+1}.
\]
The finite partition formula for the \(m\)-th derivative of
\(\e^{f(s)}\) consequently bounds its magnitude divided by \(m!\)
by \(\e X^{-m}\mathsf B_m(x)\).
Taylor's integral remainder at \(s=1\), with \(m=p+1\), has
factor \(m\int_0^1(1-s)^{m-1}\dd s=1\), proving
\eqref{eq:weighted-remainder}.
\end{proof}
The generating function has a local continuous deformation
\(A_s(q)=\exp(\log A(sq)/s)\), extended by \(A_0=1\).
It explains the grading, but the argument uses a finite Taylor
remainder. It does not exchange an infinite hierarchy with a
whole-line heat integral.

\section{The fixed-order comparison on compact time intervals}
\label{sec:comparison}
\subsection{Integrate the polynomial remainder before bounding the window}
For \(m\ge0\) define
\begin{subequations}\label{eq:gaussian-budgets}
\begin{align}
 M_m&=\sqrt2\,\e^{R+Y^2/(4\tau)+2B}
   \sum_{k=0}^m\binom{m}{k}W^{m-k}(8B)^{k/2}
                        \frac{\Gamma((k+1)/2)}{\sqrt\pi},\\
 M_m^t&=\frac{M_m}{2\tau}+\frac{M_{m+2}}{4\tau^2}.
\end{align}
\end{subequations}
Dependence on \(\tau,B,r,b\) is suppressed in these symbols.
\begin{lemma}[Complex Gaussian moments]\label{lem:moments}
For \(z\in\Omega_X(R,Y)\), \(\tau\le T\le B\), and \(q=u-X\),
\begin{subequations}\label{eq:moment-bounds}
\begin{align}
 \int_\R\abs{k_T(z-u)}\e^{\abs q}\abs q^m\dd u&\le M_m,\\
 \int_\R\abs{\partial_Tk_T(z-u)}\e^{\abs q}\abs q^m\dd u&\le M_m^t.
\end{align}
\end{subequations}
\end{lemma}
\begin{proof}
Put \(v=u-\Rea z\).
Then \(\abs q\le R+\abs v\) and
\(\max(\abs q,\abs{z-u})\le W+\abs v\).
The inequality \(\abs v\le v^2/(8T)+2T\) leaves the Gaussian
\(\e^{-v^2/(8T)}\).
Its \(k\)-th absolute moment, including the original kernel
normalization, is
\(\sqrt2(8T)^{k/2}\Gamma((k+1)/2)/\sqrt\pi\).
The binomial theorem proves the first inequality.
For the second differentiate the actual kernel:
\begin{equation}
 \partial_Tk_T(z-u)
 =\left(\frac{(z-u)^2}{4T^2}-\frac1{2T}\right)k_T(z-u).
 \label{eq:time-kernel}
\end{equation}
The same moment calculation proves \eqref{eq:gaussian-budgets},
including both inverse-time factors.
\end{proof}

At the real center \(\abs{b_{\epsilon n}}=g n^{-1/2}\).
Sectorial calibration gives \(\abs{\Rea\lambda_{\epsilon n}}\le1\)
after a fixed onset. Therefore
\begin{equation}
 \sum_{\epsilon=\pm1}\sum_{n=1}^N
       \abs{b_{\epsilon n}\e^{\lambda_{\epsilon n}q}}
 \le C_{\mathrm{pop}}gX^{1/4}\e^{\abs q},
 \qquad C_{\mathrm{pop}}=4(4\pi)^{-1/4}.
 \label{eq:population}
\end{equation}
Here \(\sum_{n\le N}n^{-1/2}\le2\sqrt N\).
Combining \cref{lem:weighted,lem:moments} with
\eqref{eq:population} gives local field and direct-time errors
\begin{equation}
 C_{\mathrm{pop}}\e\,gX^{-(p+3/4)}\mathsf B_{p+1}[M],
 \qquad
 C_{\mathrm{pop}}\e\,gX^{-(p+3/4)}\mathsf B_{p+1}[M^t].
 \label{eq:local-payments}
\end{equation}
The time bound uses moments through degree \(2p+4\).
No factor \((\log X)^{2p+2}\) appears: the powers of \(q\) were
integrated, not replaced by \(V_B\).

\subsection{Why the completion order must track the curvature order}
\label{subsec:matching}
On the window, sectorial gamma estimates give
\(G(X+q)\le C_\Gamma g\e^{\abs q}\).
If \(X-V_B\ge X/2\), \eqref{eq:rs-remainder} becomes
\begin{subequations}\label{eq:CK}
\begin{align}
 \abs{H(0,u)-f_K(u)}
 &\le C_Kg\e^{\abs q}X^{-(K/2+3/4)},\\
 C_K&=C_\Gamma\frac{2^{7/4}}7
 \Gamma((K+1)/2)1.1^{K+1}(8\pi)^{K/2+3/4}.
\end{align}
\end{subequations}
Its field and direct-time integrals cost \(C_KM_0\) and \(C_KM_0^t\)
at this power. Thus the local curvature remainder and the canonical
initial remainder have the same exponent precisely when
\begin{equation}
 K=2p.
 \label{eq:matching}
\end{equation}
This is the least integer order supporting that rate by the
displayed upper bound. It is not a lower bound on the actual
Riemann--Siegel error. Retaining only \(K=2p-1\) leaves an
\(X^{-(p+1/4)}\) allowance even though the arithmetic approximation
has improved further.

\subsection{True tails and artificial tails}
The whole-line polynomial transform in \cref{sec:hierarchy} creates
an artificial tail outside the local window. It is distinct from
the true initial tail of \eqref{eq:heat-convolution}.
Introduce
\begin{subequations}\label{eq:tail-moments}
\begin{align}
 \mathsf D_m&=\frac{\e^{Y^2/(4\tau)}}{\sqrt{4\pi\tau}}\,
      2^m(64B)^{(m+1)/2}\Gamma((m+1)/2),&
 \mathsf D_m^t&=\frac{\mathsf D_m}{2\tau}
                        +\frac{\mathsf D_{m+2}}{4\tau^2},\\
 \mathsf J_m&=\frac{\e^{Y^2/(4\tau)}}{\sqrt{4\pi\tau}}
  \sum_{k=0}^m\binom{m}{k}Y^{m-k}
                 (8B)^{(k+1)/2}\Gamma((k+1)/2),&
 \mathsf J_m^t&=\frac{\mathsf J_m}{2\tau}
                        +\frac{\mathsf J_{m+2}}{4\tau^2}.
\end{align}
\end{subequations}
When \(V_B\ge\max(2R,W+1,32B)\), for \(\abs q\ge V_B\),
\[
 -\frac{(q-\Rea d)^2}{4B}+\abs q\le-\frac{q^2}{32B},
 \qquad \abs{z-u}\le2\abs q.
\]
Splitting the last Gaussian in half and extracting its value at
\(V_B\) gives the tail moment bounds
\(\mathsf D_m\e^{-V_B^2/(64B)}\le\mathsf D_m\e^{-L^2}\)
and their direct-time versions.

The real canonical initial field has the height-sensitive estimate
\[
 \abs{H(0,X+q)}
 \le C X^{P_0}\e^{-\pi X/8+\pi\abs q/8},
 \qquad \abs q\le X/2,\qquad P_0=4.
\]
It follows from the gamma exponential and a polynomial vertical-strip
bound for zeta; the deliberately coarse exponent \(4\) is sufficient.
Only beyond \(\abs q>X/2\) do we use the global bound
\(\abs{H(0,u)}\le M_H=\int_0^\infty\abs{\Phi(v)}\dd v\).
For \(X\ge4R\), that region has
\(\abs{u-\Rea z}\ge X/4\). Splitting its Gaussian gives the
remote factor \(\e^{-X^2/(128B)}\), and division by \(g\)
costs \(\e^{\pi X/8}X^{-7/4}/g_-\).

Absorbing only fixed source constants into \(C_{\mathrm{tail}}\), set
\begin{subequations}\label{eq:tail-constants}
\begin{align}
 A_{\mathrm{mod}}&=C_{\mathrm{tail}}
                 \max(\mathsf D_0,\mathsf D_0^t),\\
 A_{\mathrm{art},p}&=C_{\mathrm{pop}}
       \max\{\mathsf H_p[\mathsf D],\mathsf H_p[\mathsf D^t]\},\\
 A_{\mathrm{far}}&=(M_H/g_-)\max(\mathsf J_0,\mathsf J_0^t).
\end{align}
\end{subequations}
The following bounds hold for both the field and the directly
computed time derivative, after division by \(g\):
\begin{center}
\begin{tabular}{@{}ll@{}}
\toprule
Contribution & Upper bound\\
\midrule
True data, \(V_B<\abs q\le X/2\) &
\(A_{\mathrm{mod}}X^{P_0-7/4}\e^{-L^2}\)\\
Artificial polynomial data, \(\abs q>V_B\) &
\(A_{\mathrm{art},p}X^{1/4}\e^{-L^2}\)\\
True data, \(\abs q>X/2\) &
\(A_{\mathrm{far}}X^{-7/4}\e^{\pi X/8-X^2/(128B)}\)\\
\bottomrule
\end{tabular}
\end{center}
No gamma function or Riemann--Siegel function is extended onto
the artificial tails.

\subsection{The theorem and its derivation}
\begin{theorem}[Fixed order on a fixed compact time family]
\label{thm:fixed}
Fix \(p\ge1\), \(r,b>0\), and \(0<\tau\le B<\infty\).
For \(K=2p\), the model \eqref{eq:model} satisfies, for all
sufficiently large \(X\),
\begin{equation}
 \normg{h-F_{X,B}^{[p,2p]}}
 \le C_{p,r,b}(\tau,B)X^{-(p+3/4)}.
 \label{eq:fixed-result}
\end{equation}
A sufficient constant is
\begin{subequations}\label{eq:fixed-constant}
\begin{align}
 A_{f,p}&=C_{2p}M_0+C_{\mathrm{pop}}\e\,\mathsf B_{p+1}[M]+3,\\
 A_{t,p}&=C_{2p}M_0^t+C_{\mathrm{pop}}\e\,\mathsf B_{p+1}[M^t]+3,\\
 C_{p,r,b}(\tau,B)&=C_\chi(48A_{f,p}+A_{t,p}).
\end{align}
\end{subequations}
The sufficient onset is specified in \cref{app:onsets}.
The frozen heat residual and its first spatial derivative vanish.
\end{theorem}
\begin{proof}
Split \eqref{eq:heat-convolution} at \(I_{X,B}\), and on that
window insert
\[
 H(0,u)=S_N(u)+Q_{X,2p}(u)+[H(0,u)-f_{2p}(u)].
\]
The \(Q\)-integral cancels exactly against the retained completion.
For each mode subtract the whole-line polynomial transform.
The remaining local difference is bounded by
\eqref{eq:local-payments}; the initial mismatch is bounded by
\eqref{eq:CK}; the three tails are bounded above.
The onset inequalities in \cref{app:onsets} put each tail below
\(X^{-(p+3/4)}\). This proves the field bound \(gA_{f,p}X^{-(p+3/4)}\)
on the enlarged chart and, independently, the time bound
\(gA_{t,p}X^{-(p+3/4)}\).

Both fields are holomorphic in \(z\). Radius-\(1/2\) Cauchy disks
fit inside the enlarged chart, and give factors \(j!2^j\le48\)
for \(0\le j\le3\). Convert from \(g\) to \(\abs{G(z)}\)
using \(C_\chi\). This proves \eqref{eq:fixed-constant}.
The residual assertion follows from \cref{prop:residual}.
\end{proof}
The theorem includes a one-point time interval \(\tau=B\), since
the underlying smooth functions are defined for \(T>0\).
Its time derivative is not a derivative of an error allowance.
The explicit constants deteriorate through
\(\e^{Y^2/(4\tau)+2B}\), powers of \(B\), and
\(\tau^{-1},\tau^{-2}\). No uniform limit at \(\tau=0\)
or \(B=\infty\) is asserted.

\subsection{Low-order formulas and the earlier first correction}
\label{subsec:low-orders}
Write \(\gamma_\epsilon=\ell_\epsilon''(X)\) and
\(\eta_\epsilon=\ell_\epsilon'''(X)\). Then
\[
 U_{\epsilon,1}=\frac{\gamma_\epsilon}{2}\calQ_T^2U_{\epsilon,0},
 \quad
 U_{\epsilon,2}=\left(\frac{\eta_\epsilon}{6}\calQ_T^3+
             \frac{\gamma_\epsilon^2}{8}\calQ_T^4\right)U_{\epsilon,0}.
\]
For one mode the first corrected expression is
\begin{equation}
 b\e^{\lambda d+T\lambda^2}
       \left[1+\frac{\gamma}{2}\{(d+2T\lambda)^2+2T\}\right].
 \label{eq:first-frozen}
\end{equation}
At \(d=0\), its bracket is
\(1+\gamma(T+2T^2\lambda^2)\).
This agrees with the arithmetic value of the earlier
target-dependent correction \cite{ans10},
\[
 h_{\epsilon n}(z)\e^{T\beta_{\epsilon n}(z)^2}
    [1+\gamma_\epsilon(z)(T+2T^2\beta_{\epsilon n}(z)^2)],
 \quad
 \beta_{\epsilon n}(z)=\partial_z\log h_{\epsilon n}(z),
\]
evaluated at \(z=X\). In that formula, however, the gamma data
vary with the target. Its spatial derivatives differ from those
of \eqref{eq:first-frozen}, and its residual is
\(O(\abs G X^{-7/4})\), not identically zero.
Its \(K=1\) completion gives a total \(X^{-5/4}\) error allowance.
For the frozen model, the local first-level allowance is
\(X^{-7/4}\); retaining \(K=2\) supports it. At the second level
the mixed curvature-square term and \(C_0,\ldots,C_4\) must all
remain to support \(X^{-11/4}\).

Thus higher approximation means retaining all partitions of the
chosen fluctuation weight and matching the completion order.
Adding only one higher derivative of gamma does not implement
the next level.

\section{Growing order and outer-envelope accuracy}
\label{sec:growing}
We now specialize to \(1\le T\le2\) and
\(V=8\sqrt2L\). The fixed-order theorem cannot be used by
substituting \(p=p(X)\) into an unspecified \(C_p\).
We instead retain the pre-absorption error estimates and bound
their constants simultaneously.

\subsection{Factorial bounds for the actual finite constants}
Put
\begin{align}
 E_*&=\max\{1,\sqrt2\e^{R+Y^2/4+4},
                         4\sqrt2\e^{Y^2/4}\},\\
 S_*&=\max\{1,W+4,16\sqrt2\},\qquad C_*=\max(1,A,D).
\end{align}
The half-integer gamma recurrence gives
\(\Gamma((m+1)/2)/\sqrt\pi\le m!\).
The finite moment sums consequently imply
\[
 M_m,\mathsf D_m\le E_*m!S_*^m,\qquad
 M_m^t,\mathsf D_m^t\le E_*(m+2)!S_*^{m+2}.
\]
Grouping the sums in \eqref{eq:majorant-polynomials} gives
\begin{subequations}\label{eq:factorial-majorants}
\begin{align}
 \mathsf B_m[M]&\le E_*(8C_*S_*^2)^m(2m)!,\\
 \mathsf B_m[M^t]&\le E_*S_*^2(8C_*S_*^2)^m(2m+2)!,\\
 \max\{\mathsf H_p[\mathsf D],\mathsf H_p[\mathsf D^t]\}
 &\le E_*S_*^2(p+1)(2C_*S_*^2)^p(2p+2)!.
\end{align}
\end{subequations}
For example, in the first line bound
\((m+k)!\le(2m)!\), \(S_*^{m+k}\le S_*^{2m}\),
and \(\sum_{k=1}^m\binom{m-1}{k-1}/k!\le2^{m-1}\).
Also
\begin{equation}
 C_{2p}\le C_RQ_R^p p!,\quad
 C_R=C_\Gamma\frac{2^{7/4}}7\sqrt\pi\,1.1(8\pi)^{3/4},
 \quad Q_R=1.1^2(8\pi).
 \label{eq:rs-factorial}
\end{equation}
Let \(Q_*=\max(1,8C_*S_*^2,2C_*S_*^2,Q_R)\),
\(T_*=\max(1,E_*S_*^2,C_R)\). The common explicit majorant
\[
 \mathcal J(p)=T_*(p+2)Q_*^{p+1}(2p+4)!
\]
bounds each cost in \eqref{eq:factorial-majorants} and \(C_{2p}\).
In particular
\begin{equation}
 \log\mathcal J(p)=O((p+1)\log(p+2)).
 \label{eq:log-growth}
\end{equation}
This bound precedes any choice \(p=O(L)\).
It follows directly from \(n!\le n^n\), with constants fixed
by the chart. The original finite sums remain the operative
constants when checking a particular height.

\subsection{Keep the tails before normalizing}
Write
\[
 Q_p=48\{C_{2p}M_0+C_{\mathrm{pop}}\e\,\mathsf B_{p+1}[M]\}
       +C_{2p}M_0^t+C_{\mathrm{pop}}\e\,\mathsf B_{p+1}[M^t].
\]
The proof of \cref{thm:fixed}, before tail absorption, gives
\begin{align}
 \lVert h-F_X^{[p,2p]}\rVert_{\gamma;X;1,2}^{(3,1)}
 \le C_\chi\bigl\{&
 Q_pX^{-(p+3/4)}
 +49(A_{\mathrm{mod}}X^{P_0-7/4}
                   +A_{\mathrm{art},p}X^{1/4})\e^{-L^2}\nonumber\\
 &+49A_{\mathrm{far}}X^{-7/4}
                       \e^{\pi X/8-X^2/256}\bigr\}.
 \label{eq:unabsorbed}
\end{align}
This estimate has common structural onsets for \(1\le p\le cL\),
each fixed \(c>0\): replace the retained-order condition by
\(2cL\le(X-V)/(4\pi)\). All other structural gates are
independent of \(p\).

The old fixed-order absorption condition
\[
 L^2\ge(p+1)\log X+\Lp A_{\mathrm{art},p}
\]
fails already at \(p=\lfloor L\rfloor\), because
\((p+1)\log X>L^2\).
That failure concerns an unnecessarily strong intermediate bound.
The unabsorbed estimate \eqref{eq:unabsorbed} remains available.

\subsection{A stronger Gaussian tail with visible costs}
Fix \(0<\theta<1\), put \(\delta=(1-\theta)/2\) and \(t_R=1+R/4\).
For \(\abs q\ge V\),
\[
 -\frac{(q-\Rea d)^2}{8}+\abs q
 \le-\frac{q^2}{8}+t_R\abs q
 \le-\frac{\theta q^2}{8}-\frac{\delta q^2}{8}
                         +\frac{2t_R^2}{\delta}.
\]
Define
\begin{subequations}\label{eq:theta-moments}
\begin{align}
 \mathsf D_m^\theta
 &=\frac{\e^{Y^2/4+2t_R^2/\delta}}{\sqrt{4\pi}}\,
          2^m(8/\delta)^{(m+1)/2}\Gamma((m+1)/2),\\
 (\mathsf D_m^\theta)^t
 &=\mathsf D_m^\theta/2+\mathsf D_{m+2}^\theta/4.
\end{align}
\end{subequations}
Extracting \(\e^{-\theta V^2/8}=\e^{-16\theta L^2}\)
and integrating the remaining Gaussian proves the refined tail
bound. Define \(A_{\mathrm{mod}}^\theta,A_{\mathrm{art},p}^\theta\)
by \eqref{eq:tail-constants} with these moments.
Then \eqref{eq:unabsorbed} holds with these constants and
\(\e^{-16\theta L^2}\) replacing \(\e^{-L^2}\).

The cost of the sharper tail is explicit. Use
\[
 E_\theta=\max\{1,\sqrt2\e^{R+Y^2/4+4},
        \e^{Y^2/4+2t_R^2/\delta}\sqrt{2/\delta}\},
 \quad
 S_\theta=\max\{1,W+4,4\sqrt{2/\delta}\}
\]
in \eqref{eq:factorial-majorants}.
Thus \(\log^+A_{\mathrm{art},p}^\theta
=O_\theta((p+1)\log(p+2))\).
The parameter \(\theta\) is fixed before the height tends to
infinity; the constants are not uniform as \(\theta\to1\).

\begin{theorem}[Controlled growing order]\label{thm:growing}
Fix \(r,b>0\), \(c>0\), and \(a\) with \(1/8<a\le c\).
For \(p=\lfloor aL\rfloor\), \(K=2p\), and every
\[
 0<\eta<\min(a-1/8,127/8),
\]
the full frozen model on \(1\le T\le2\) satisfies, eventually,
\begin{equation}
 \normo{h-F_X^{[p,2p]}}\le\e^{-\eta L^2}.
 \label{eq:growing-result}
\end{equation}
In particular \(p=\lfloor L/2\rfloor\), \(K=2p\) gives
\(\normo{h-F_X^{[p,2p]}}\le\e^{-L^2/4}\).
The estimate retains both arms, every \(n\le N\), and the bounded
completion integral.
\end{theorem}
\begin{proof}
Choose a fixed \(\theta<1\) such that
\(\eta<16\theta-1/8\).
By \eqref{eq:log-growth}, the logarithm of the source/local
constant is \(O_a(L\log L)\).
After \eqref{eq:outer-conversion}, its exponent is at most
\[
 -(p+3/4)\log X+L^2/8+O_a(L\log L)
 =-(a-1/8)L^2+O_a(L\log L).
\]
The moderate and artificial tails have exponents at most
\[
 -(16\theta-1/8)L^2+O_{a,\theta}(L\log L).
\]
The remote true tail is smaller still, with a negative quadratic
term in \(X\).
The strict inequalities for \(\eta\) make these bounds summably
smaller than the target after a common onset.
The finite four-part allocation is given in \cref{app:growing-gates}.
Taking \(\theta=1/2\), \(a=1/2\), \(\eta=1/4\) proves the
stated concrete choice.
\end{proof}
The range is sufficient, not asserted sharp. A cap \(p\le cL\)
alone would include \(p=1\) and would not imply outer decay.
The bound is relative to the positive envelope, including near
zeros of the signed field. Selection of \(p(X)\) does not permit
setting \(X=\Rea z\) during differentiation or differentiating
the floor. Each chart remains frozen.

The earlier ansatz investigation already supplied a separate
growing-degree outer approximation \cite{ans05}.
The present result establishes the corresponding accuracy for
the fluctuation grading, with simultaneous factorial costs and
the same canonical completion. It does not claim that the first
outer approximation has only now been obtained.

\section{Analytic regularity and quadrature for the retained completion}
\label{sec:quadrature}
The approximation theorem leaves a finite bounded integral.
Controlling its evaluation requires more than the real size
estimate in \cref{subsec:rs-input}. We now specialize throughout
to \(1\le T\le2\), \(p=\lfloor L/2\rfloor\), \(K=2p\).
The statements below concern the finite retained expression, not
an analytic continuation of the real Riemann--Siegel error theorem.

\subsection{Remove apparent poles before estimating coefficients}
The exact coefficient identities used here are those recorded in
Arias de Reyna, Part II, (29), (62), (86)--(89)
\cite{ariasII}, attributed there to the published Part I
\cite{ariasI}. The underlying Part I proof has not been independently
rechecked in this project. Part II contains preliminary precision
arguments; its unproved technical inequalities are not imported.
We use the named identities and derive the following bounds directly.

In the same complex normalization as \cref{subsec:rs-input}, set
\begin{subequations}\label{eq:rs-coefficient-identity}
\begin{align}
 F_0(\chi)&=
 \frac{\e^{\pi\ii(\chi^2/2+3/8)}
                     -\ii\sqrt2\cos(\pi\chi/2)}{2\cos(\pi\chi)},\\
 C_k(\chi,1/2)&=\pi^{-2k}
  \sum_{j=0}^{\lfloor3k/2\rfloor}
       \left(\frac{\pi}{2\ii}\right)^j
                  d_j^{(k)}F_0^{(3k-2j)}(\chi).
\end{align}
\end{subequations}
At \(\chi=n+1/2\) the two numerator terms are both
\(\ii(-1)^{n(n+1)/2}\). Every zero of the denominator is simple
and removable, so \(F_0\) is entire.
This fact alone would not make direct numerical division near
those points stable.

\begin{lemma}[Weighted coefficient recurrence]\label{lem:weighted-rs}
Let \(e_j^{(k)}=(3k-2j)!\,d_j^{(k)}\), with \(e_0^{(0)}=1\)
and out-of-range entries zero. For \(m=3k-2j>0\),
\begin{equation}
 e_j^{(k)}
 =\frac{(m-1)(m-2)}4e_j^{(k-1)}-e_{j-2}^{(k-1)}.
 \label{eq:weighted-rs}
\end{equation}
The first term is structurally absent when \(m=1,2\).
For positive even \(k\), with \(j=3k/2\), the endpoint is
\begin{equation}
 e_j^{(k)}=-\sum_{h=1}^{j}\frac{(-1)^h}{h!}e_{j-h}^{(k)}.
 \label{eq:rs-endpoint}
\end{equation}
The values and their positive recurrence majorants satisfy
\begin{equation}
 \max_j\abs{e_j^{(k)}}\le E_k:=8^k(k!)^2.
 \label{eq:weighted-rs-bound}
\end{equation}
\end{lemma}
\begin{proof}
Multiply the exact coefficient recurrence at \(\sigma=1/2\)
by \(m!\). The previous-order derivative indices are \(m-3\)
and \(m+1\), which gives \eqref{eq:weighted-rs}.
The exact endpoint relation gives \eqref{eq:rs-endpoint}.
The sum of absolute interior coefficients is at most
\((9/4)k^2\). The endpoint contributes a factor at most
\(\sum_{h\ge1}1/h!<2\) times the largest interior majorant.
Thus a row costs at most \((9/2)k^2\) times the preceding row.
Induction with \(8k^2\) proves \eqref{eq:weighted-rs-bound},
also for intermediate positive sums.
\end{proof}

\begin{proposition}[Explicit complex coefficient bounds]
\label{prop:rs-complex}
Put
\[
 B_F=\tfrac12\{\e^{2\pi}+\sqrt2\cosh(\pi/2)\}.
\]
For \(\abs{\Rea\chi}\le3/2\), \(\abs{\Ima\chi}\le1/2\),
\begin{equation}
 \abs{C_k(\chi,1/2)}\le
 D_k:=\frac{B_F}{1-\pi/8}
                   (64/\pi^2)^k(k!)^2.
 \label{eq:complex-Ck}
\end{equation}
At real \(\chi\in[-1,1]\), for each fixed \(s\ge0\),
\begin{equation}
 \abs{\partial_\chi^sC_k(\chi,1/2)}
 \le B_F\pi^{-2k}E_k
       \sum_{j=0}^{\lfloor3k/2\rfloor}
       (\pi/2)^j\frac{(3k-2j+s)!}{(3k-2j)!}.
 \label{eq:real-Ck-derivative}
\end{equation}
\end{proposition}
\begin{proof}
On the rectangle \(\abs{\Rea\chi}\le2,\abs{\Ima\chi}\le1\),
the quotient is bounded by \(B_F\) on its boundary:
the denominator has modulus at least \(2\) on the vertical
sides and at least \(2\sinh\pi\) on the horizontal sides.
Maximum modulus applies to the entire continuation.
Radius-\(1/2\) Cauchy disks then give
\(\abs{F_0^{(m)}}\le B_Fm!2^m\) on the smaller rectangle.
Substitute \eqref{eq:weighted-rs-bound} into
\eqref{eq:rs-coefficient-identity} and sum
\(\sum_{j\ge0}(\pi/8)^j\), proving \eqref{eq:complex-Ck}.
At real \(\chi\in[-1,1]\), radius-one Cauchy disks prove
\eqref{eq:real-Ck-derivative}.
\end{proof}
In particular \(\log^+D_k=O((k+1)\log(k+2))\).
The complex normalization is essential; the real cosine quotient
\(\mathcal C\) in \eqref{eq:short} is not substituted for \(F_0\).

\subsection{Holomorphic extension of each finite branch}
Split \(I_X\) at its exact thresholds \(4\pi n^2\).
On a branch \(J\), fix \(m\) and use
\(\chi_m(u)=1-2(\sqrt{u/(4\pi)}-m)\).
If \(A_m(u)\) denotes the finite complex expression inside
the real part in \eqref{eq:real-completion}, then
\[
 A_m(u)+\overline{A_m(\bar u)}
\]
is holomorphic wherever the displayed finite factors are, and
agrees with the real initial correction on \(J\).
Add the two analytic transfer arms. This extends the retained
branch \(Q_m\) holomorphically; it makes no assertion about the
discarded remainder off the real axis.

Use the Bernstein ellipse of parameter \(\rho=2\) for each \(J\).
It lies within \(\abs{u-X}\le2V\), \(\abs{\Ima u}\le3V/4\),
and every point on it is within \(V\) of a real point of \(J\).
The gates
\[
 2V\le X/4,\qquad V/\sqrt{2\pi X}\le1/2
\]
ensure that the branch stays in the coefficient rectangle of
\cref{prop:rs-complex}. Write
\begin{subequations}\label{eq:ellipse-data}
\begin{align}
 a_*&=\sqrt{(X-2V)/(4\pi)},\\
 H_M&=\exp(2V\abs{\ell_+'(X)}+2HV^2/X),&
 H_U&=\exp(V(L+1)/2),\\
 \mathcal T_X&=\sum_{n\in\mathcal T}n^{-1/2}\e^{V\log n},&
 B_Q&=2H_M\left(\mathcal T_X+
                    H_Ua_*^{-1/2}\sum_{k=0}^KD_ka_*^{-k}\right).
\end{align}
\end{subequations}
Here \(\mathcal T\) is the union of actual transfer indices
over the branches; it contains at most one index.
The gamma bound follows from \cref{lem:coeff}.
For \(U\), integrate its logarithmic derivative
\(-\ii\Log(u/(4\pi))/4\), imposing the eventual gate
\(\sup_{\abs{q}\le2V}\abs{\Log((X+q)/(4\pi))}\le L+1\).
Thus \(\abs{Q_m(u)/g}\le B_Q\) on every branch ellipse.

Define nonnegative kernel polynomials
\[
 \mathsf K_a(v)=a!\sum_{j=0}^{\lfloor a/2\rfloor}
       \frac{(v/2)^{a-2j}(1/4)^j}{(a-2j)!j!},
 \qquad 0\le a\le3,
\]
and
\begin{equation}
 K_E=\frac{\e^{(Y+3V/4)^2/4}}{\sqrt{4\pi}}
       \left[\max_{0\le a\le3}\mathsf K_a(W+2V)
                       +\mathsf K_2(W+2V)\right],
 \qquad B_E=B_QK_E.
 \label{eq:ellipse-kernel}
\end{equation}
The sum of the maximum spatial channel and the time channel
of the normalized integrand is bounded by \(B_E\).
For the latter use \(\partial_Tk_T=\partial_z^2k_T\);
these are actual kernel derivatives.
The estimates retain the off-real kernel growth, giving
\(\log^+B_E=O(L^2)+O(K\log(K+2))\).

\subsection{Analytic Gauss error}
\begin{lemma}[A complex-valued Gauss bound]\label{lem:gauss}
If an analytic integrand is bounded by \(B\) on the
Bernstein \(\rho\)-ellipse of a real interval \(J\), its
\(n\)-node Gauss--Legendre rule has error at most
\[
 4\abs J B\,\frac{\rho^{-2n+1}}{\rho-1}.
\]
The statement also holds for a finite-dimensional vector integrand
in any norm.
\end{lemma}
\begin{proof}
Under the Joukowski map, Cauchy's coefficient estimate gives
Chebyshev coefficients bounded by \(2B\rho^{-j}\).
Truncation at degree \(2n-1\) has uniform remainder at most
\(2B\rho^{-2n+1}/(\rho-1)\).
Gauss weights are positive, sum to \(\abs J\), and the rule
is exact through degree \(2n-1\) \cite{dlmfgauss}.
Subtract the polynomial; the integral and the quadrature
each cost at most \(\abs J\) times the uniform remainder.
The same argument applies to the channel norm. It uses no
real mean-value formula for complex-valued data.
\end{proof}
Summing branch lengths, \(\sum_J\abs J=2V\), and setting
\(\rho=2\) gives
\begin{equation}
 E_{\mathrm{quad}}\le16VB_E4^{-n}.
 \label{eq:quadrature-bound}
\end{equation}
For any assigned center-normalized budget \(\varepsilon/20\),
it is sufficient to take
\begin{equation}
 n\ge\max\left\{1,\left\lceil
           \frac{\log(320VB_E/\varepsilon)}{\log4}\right\rceil\right\}.
 \label{eq:quadrature-order}
\end{equation}
When \(-\log\varepsilon=O(L^2)\) and \(K=O(L)\), this is
\(n=O(L^2)\) on each of at most two branches.
Node, weight, and input-data errors have not been included
in \eqref{eq:quadrature-bound}; they are paid next.

\section{Evaluation error budgets}
\label{sec:evaluation}
The theorem in this section is a reusable implication from delivered-data
contracts to field accuracy. Sections~\ref{sec:providers}--\ref{sec:effective}
subsequently construct those data and prove an effective total result.

\subsection{Normalize first and allocate the total error}
Take \(p=\lfloor L/2\rfloor\), \(K=2p\).
Increase the strict-margin onset in \cref{thm:growing} so that
the mathematical truncation uses at most
\(\tfrac12\e^{-L^2/4}\) of the outer budget.
By \eqref{eq:outer-conversion} a sufficient center-normalized
evaluation tolerance is
\begin{equation}
 \varepsilon_X=\frac{c_{\mathrm o}}{2C_\chi}\e^{-3L^2/8}.
 \label{eq:epsilonX}
\end{equation}
Write \(\beta_{\epsilon n}=b_{\epsilon n}/g\) and
\(P_{\epsilon,p}(q)=\sum_{m=0}^{2p}t_{\epsilon m}q^m\).
The normalized exact expression is
\begin{align}
 F/g={}&\sum_{\epsilon,n\le N}\beta_{\epsilon n}
       \e^{\lambda_{\epsilon n}d+T\lambda_{\epsilon n}^2}
       \sum_{m=0}^{2p}t_{\epsilon m}
              \calH_m(d+2T\lambda_{\epsilon n},T)\nonumber\\
 &+\int_{X-V}^{X+V}k_T(z-u)Q_{X,K}(u)/g\dd u.
 \label{eq:normalized-evaluation}
\end{align}
Let \(A_0\) bound its combined center-scaled channel magnitudes.
If \(\abs{\widehat g/g-1}\le\varepsilon_g\le1\), it suffices that
\begin{equation}
 \varepsilon_g A_0\le\varepsilon_X/2,\qquad
 E_{\mathrm{ar}},E_{\mathrm{RS}},E_{\mathrm{quad}},
 E_{\mathrm{geo}},E_{\mathrm{round}}\le\varepsilon_X/20.
 \label{eq:allocation}
\end{equation}
Indeed the center-normalized output error is bounded by
\(\varepsilon_gA_0+(1+\varepsilon_g)\sum E_i\le\varepsilon_X\).
The five categories are, respectively, arithmetic inputs and
preprocessing, retained-integral data, exact quadrature, geometry,
and remaining operation rounding. Scale error is separate.

\subsection{Arithmetic input conditioning}
Use the moments at \(\tau=1,B=2\), and set
\[
 \mathfrak M_m=48M_m+M_m^t,\quad
 \mathcal D_p=\mathsf H_p[\mathfrak M],\quad
 \mathcal D_p^{(1)}=\mathsf H_p[q\mathfrak M],\quad
 \mathcal S_p=\sum_{m=0}^{2p}\mathfrak M_m.
\]
The notation \(\mathsf H_p[q\mathfrak M]\) pairs its \(q^m\)
coefficient with \(\mathfrak M_{m+1}\).
The factor \(48\) includes the buffered spatial Cauchy costs.

\begin{lemma}[Finite coefficient sensitivity]\label{lem:sensitivity}
For \(1\le r\le p\),
\begin{equation}
 \partial_{c_r}P_p(q)=q^{r+1}P_{p-r}(q).
 \label{eq:coefficient-sensitivity}
\end{equation}
For perturbations \(\abs{\delta c_r}\le\varepsilon_c\le1\),
let \(\mathsf C_j^+\) denote \(\mathsf C_j\) with \(A\) replaced
by \(A+1\), and \(\mathsf C_0^+=1\). Then a sufficient condition
sum is
\begin{equation}
 \mathcal C_p=\sum_{r=1}^p\sum_{j=0}^{p-r}
                     \mathsf C_j^+[q^{r+1}\mathfrak M].
 \label{eq:condition-sum}
\end{equation}
\end{lemma}
\begin{proof}
Differentiate the local generating function: differentiation
in \(c_r\) multiplies it by \(q^{r+1}s^r\), proving
\eqref{eq:coefficient-sensitivity} coefficient by coefficient.
Along the segment of perturbed coefficients,
\(\abs{c_r'}\le(A+1)D^{r-1}\), after dropping the favorable
\(X^{-r}\). Apply the positive partition majorant and integrate
the derivative along that segment. Then use the Gaussian
channel moments, giving \eqref{eq:condition-sum}.
\end{proof}

Impose
\(D_0/X\le1/2-\pi/8\) and \(\varepsilon_\lambda\le1/2\).
The sectorial slope formula ensures that the original slopes
and their perturbed connecting segments have real part of
absolute value at most one.
For
\[
 \abs{\delta\beta_{\epsilon n}}\le\varepsilon_\beta n^{-1/2},
 \quad\abs{\delta\lambda_{\epsilon n}}\le\varepsilon_\lambda,
 \quad\abs{\delta c_{\epsilon r}}\le\varepsilon_c,\quad
 \abs{\delta t_{\epsilon m}^{\mathrm{store}}}\le\varepsilon_t,
\]
one has
\begin{equation}
\begin{split}
 E_{\mathrm{ar}}\le C_{\mathrm{pop}}X^{1/4}
 \bigl[\varepsilon_\beta\mathcal D_p
 +(1+\varepsilon_\beta)
 \{\varepsilon_\lambda\mathcal D_p^{(1)}
               +\varepsilon_c\mathcal C_p+\varepsilon_t\mathcal S_p\}\bigr].
\end{split}
\label{eq:arithmetic-error}
\end{equation}
To verify the estimate, change \(\beta\) first with the exact
polynomial and slope, then the slope, then the shape coefficients
and stored output. Differentiation in \(\lambda\) inside the
real Gaussian transform inserts \(q\).
The real tilt remains bounded by \(\e^{\abs q}\), so the same
moments apply. This ordering avoids using an unperturbed
polynomial majorant for perturbed coefficients.
The error \(\varepsilon_t\) includes the certified error of
preprocessing at the supplied \(c_r\); it is charged here
only once.

The logarithms of all finite condition sums are
\(O((p+1)\log(p+2))\), by the factorial bounds and the finite
partition formulas. There is no division by a possibly vanishing
field or correction.

\subsection{Large phases and very small scales}
At real \(X\),
\[
 \beta_{\epsilon n}=n^{-1/2}
    \e^{\ii(\theta_\epsilon(X)-\epsilon X\log n/2)},\qquad
 g=\e^{\Rea\ell_+(X)}.
\]
If the error in the carrier logarithm is at most \(q_{\mathrm{err}}\),
then
\(\abs{\delta\beta_{\epsilon n}}
\le n^{-1/2}\e^{q_{\mathrm{err}}}q_{\mathrm{err}}\).
A sufficient allowance is
\begin{align}
 q_{\mathrm{err}}\le{}&
 \abs{\delta\theta}
 +(X/2+1/2)\abs{\delta\log n}
 +\tfrac12\log n\abs{\delta X}
 +\tfrac12\abs{\delta X\,\delta\log n}\nonumber\\
 &+\abs{\delta(\Rea\log M-\log g)}.
 \label{eq:phase-budget}
\end{align}
Also
\(\abs{\delta\lambda}\le\abs{\delta\ell'}+
\abs{\delta\log n}/2\) at fixed \(X\).
Thus logarithms of mode indices must be accurate to order
\(\varepsilon_\beta/X\).
Unreduced phases have size \(O(XL)\), requiring
\(O(L+\log L)\) guard bits when formed by relative arithmetic.
Normalized gamma ratios should be formed from log differences.
Dividing separately underflowed values of \(M\) and \(g\)
does not implement this specification.
Sufficient exponent range is part of the model.

\subsection{Real-node data and geometry}
At a real node on branch \(m\), write
\[
 A=(-1)^{m-1}M(w_+(u))U(u)a(u)^{-1/2}/g,\qquad R_k=a(u)^{-k}.
\]
If the errors in \(A,R_k,C_k\), and the normalized transfer
are bounded by \(\varepsilon_A,\varepsilon_{R_k},
\varepsilon_{C_k},\varepsilon_{\mathrm{tr}}\), respectively, then
\begin{equation}
\begin{split}
 \abs{\widehat Q_m/g-Q_m/g}
 \le\varepsilon_{\mathrm{tr}}+2\sum_{k=0}^K
 \bigl[&\varepsilon_A R_k\abs{C_k}\\
 &+(\abs A+\varepsilon_A)
 \{\varepsilon_{R_k}\abs{C_k}
           +(R_k+\varepsilon_{R_k})\varepsilon_{C_k}\}\bigr].
\end{split}
\label{eq:node-data-error}
\end{equation}
This is just the product rule for finite perturbations.
It retains every coefficient and the exact transfer.

Geometry errors use a local strip bound rather than the much
larger branch-ellipse bound. Require \(V\ge3\),
\[
 4(V+2)/\sqrt{2\pi X}\le1/2,\qquad 2(V+2)\le X/4.
\]
Define
\begin{subequations}\label{eq:local-strip}
\begin{align}
 H_M^{\mathrm{loc}}
 &=\exp\left\{\frac{\pi}{8}(V+2)+\frac L4+
             \frac{D_0(V+3)}X+\frac{H(V+3)^2}{2X}\right\},\\
 H_U^{\mathrm{loc}}&=\e^{(L+1)/4},\\
 B_Q^{\mathrm{loc}}
 &=2H_M^{\mathrm{loc}}\left[
       \sum_{n\in\mathcal T} n^{-1/2}n^{1/2}
       +H_U^{\mathrm{loc}}a_*^{-1/2}\sum_{k=0}^KD_ka_*^{-k}\right],\\
 K_{\mathrm{loc}}
 &=\frac{\e^{(Y+1)^2/4}}{\sqrt{4\pi}}
       \left[\max_{a\le3}\mathsf K_a(W+V+4)
                         +\mathsf K_2(W+V+4)\right],\\
 B_{\mathrm{loc}}&=B_Q^{\mathrm{loc}}K_{\mathrm{loc}}.
\end{align}
\end{subequations}
On \(\abs{\Rea u-X}\le V+2,\abs{\Ima u}\le1\), these bound
the normalized integrand channels. For gamma, use
\(\ell_+'(X)=-\pi/8+\ii L/4+O(D_0/X)\) and bound real and
imaginary displacements separately. This proves the displayed
\(H_M^{\mathrm{loc}}\); using \(\abs{\ell'}\) times the full
window instead would lose the useful bound
\[
 \log^+B_{\mathrm{loc}}=O(L)+O(K\log(K+2)).
 \label{eq:local-strip-growth}
\]
Radius-\(1/2\) Cauchy estimates bound each derivative in \(u\)
by twice its local bound.

Let \(\varepsilon_w\) be the sum of absolute physical weight
errors, and let real approximate nodes have maximum displacement
\(\varepsilon_u\le1/4\). Use their fixed reference branch labels.
If omitted boundary strips have total length at most
\(6\varepsilon_b\), \(0\le\varepsilon_b\le1/4\), then
\begin{equation}
 E_{\mathrm{geo}}\le B_{\mathrm{loc}}
   \{\varepsilon_w+2(2V+\varepsilon_w)\varepsilon_u
                                      +6\varepsilon_b\}.
 \label{eq:geometry-error}
\end{equation}
The terms follow by changing weights, then nodes, and integrating
over omitted strips. Such strips must lie inside the true fixed
branches. There are at most three boundaries, so two-sided
width-\(\varepsilon_b\) strips have the stated total length.
For exact reference endpoints take \(\varepsilon_b=0\);
endpoint errors propagated into node and weight errors are not
charged again. Rules extending outside the original window need
a new length and ellipse analysis.

If \eqref{eq:node-data-error} is bounded by \(\varepsilon_Q\)
at every approximate node, then the separate data cost is
\begin{equation}
 E_{\mathrm{RS}}\le(2V+\varepsilon_w)K_{\mathrm{loc}}\varepsilon_Q.
 \label{eq:RS-data-budget}
\end{equation}
An explicit scale bound for \eqref{eq:allocation} is
\begin{equation}
 A_0=C_{\mathrm{pop}}X^{1/4}\mathsf H_p[\mathfrak M]
                                      +2VB_{\mathrm{loc}}.
 \label{eq:A0}
\end{equation}
Its logarithm is \(O(L\log L)\).

\subsection{A coefficient construction away from apparent poles}
Let \(h_m(\chi)=F_0^{(m)}(\chi)/m!\).
On the rectangle \(\Gamma\) with vertices \(\pm2\pm\ii\),
Cauchy's formula gives
\begin{equation}
 h_m(\chi)=\frac1{2\pi\ii}\oint_\Gamma
                         \frac{F_0(\zeta)}{(\zeta-\chi)^{m+1}}\dd\zeta.
 \label{eq:contour-jet}
\end{equation}
For \(\abs{\Rea\chi}\le3/2,\abs{\Ima\chi}\le1/2\), use
a Gauss rule with \(q_c\) nodes on each straight edge and ellipse
parameter \(\rho_c=9/8\).
The edge ellipses have semiaxis ratios \(145/144,17/144\);
they satisfy
\[
 \abs{\Rea\zeta}\le9/4,\quad\abs{\Ima\zeta}\le5/4,\quad
 \abs{2\cos(\pi\zeta)}\ge1,\quad
 \operatorname{dist}(\zeta,\chi)\ge1/4.
\]
For the vertical edges the real deviation from \(\pm2\) is
at most \(17/144\); for the horizontal edges the imaginary
distance from the real line is at least \(1-34/144\).
These give the denominator bounds and the stated separation.
With \(B_c=\e^{45\pi/16}+\sqrt2\cosh(5\pi/8)\), the integrand
is bounded by \(B_c4^{m+1}\).
The total contour length is \(12\). Applying \cref{lem:gauss}
and the factor \(1/(2\pi)\) proves
\begin{equation}
 \varepsilon_{h,m}^{\mathrm{quad}}
 \le\frac{24}{\pi}B_c4^{m+1}
                  \frac{\rho_c^{-2q_c+1}}{\rho_c-1}.
 \label{eq:contour-error}
\end{equation}
For all \(m\le3K\), a requested error
\(\exp[-3L^2/8-O(L\log L)]\) is reached with
\(q_c=O(L^2+L\log L)\).
The identities now take the normalized form
\[
 C_k=\pi^{-2k}\sum_j(\pi/(2\ii))^j e_j^{(k)}h_{3k-2j}.
\]
Each perturbation term is bounded by
\(\abs{e_j}\varepsilon_h+\abs{h_{3k-2j}}\varepsilon_e+
\varepsilon_e\varepsilon_h\), times the displayed powers of \(\pi\).
Errors in those powers also require their product bounds.
This construction supplies an analytical way around the apparent
poles; primitive, node, weight, and rounding errors still have to
be certified. The fixed contour samples can be reused for all
outer quadrature nodes.

\subsection{Rounding and sufficient requested precision}
Use real-component arithmetic with unit roundoff \(u\), sufficient
exponent range, and certified bounds \(\abs a\le A_a\),
\(\abs b\le A_b\), with propagated errors \(E_a,E_b\).
For addition, multiplication, and division by an exact positive
integer \(j\), valid updates are
\begin{subequations}\label{eq:roundoff-updates}
\begin{align}
 E_+&=E_a+E_b+u(A_a+A_b+E_a+E_b),\\
 E_\times&=A_aE_b+A_bE_a+E_aE_b
                          +u(A_a+E_a)(A_b+E_b),\\
 E_{\div j}&=[E_a+u(A_a+E_a)]/j.
\end{align}
\end{subequations}
Noninteger reciprocals require certified separation and their
own error. Nonlinear primitives are not automatically covered
by these arithmetic updates. For example, if real exponential
relative error and sine/cosine absolute error are at most \(u\),
then, for \(u\le1/16\), a complex exponential with exact real
part at most \(C_0\) and argument error \(E_{\mathrm{arg}}\) has
error at most
\[
 \e^{C_0}(\e^{E_{\mathrm{arg}}}-1)
                  +4u\e^{C_0+E_{\mathrm{arg}}}.
\]

The exact derivative channels are also finite recurrences:
\begin{subequations}\label{eq:channel-formulas}
\begin{align}
 \partial_d^a(\beta\e^{\lambda d+T\lambda^2}\calH_m)
 &=\beta\e^{\lambda d+T\lambda^2}
   \sum_{j\le\min(a,m)}\binom aj(m)_j
                  \lambda^{a-j}\calH_{m-j},\\
 \partial_T(\beta\e^{\lambda d+T\lambda^2}\calH_m)
 &=\beta\e^{\lambda d+T\lambda^2}
   [\lambda^2\calH_m+2m\lambda\calH_{m-1}
                           +m(m-1)\calH_{m-2}].
\end{align}
\end{subequations}
All \(\calH\)'s here have argument \(d+2T\lambda,T\), and
\((m)_j=m!/(m-j)!\).
Evaluation needs retained moments only through \(2p\).
The error analysis still needs moments through \(2p+4\).

To bound intermediate magnitudes, take
\(\Lambda=1+L/4\ge\abs\lambda\) after an additional slope gate,
\(t_*=W+4\Lambda\), and
\[
 h_0^*=1,\quad h_1^*=t_*,\quad
 h_m^*=t_*h_{m-1}^*+4(m-1)h_{m-2}^*.
\]
Completing the square in \(\Ima\lambda\), with
\(\abs{\Rea\lambda}\le1\), gives
\[
 \abs{\e^{\lambda d+T\lambda^2}}
                  \le\e^{R+2+Y^2/4}.
\]
Thus no artificial \(\e^{+L^2}\) carrier bound is required
in the actual real-node circuits.
The triangular coefficient recurrence, moment recurrence,
weighted Riemann--Siegel recurrence, and finite linear sums
have positive path majorants of size
\((N+1)\exp(O(L\log L))\) times a polynomial in the node counts.
One can see this by expanding the triangular recurrences:
their depth is \(O(p+K)\), their factors are polynomial in
\(p,K,L\), and their partition paths have
\(\exp(O((p+K)\log(p+K+2)))\) total weight.
Linear accumulation contributes its actual length, not
an exponential in \(N\). Impose
\(uS_{\mathrm{lin}}\le1/4\), where \(S_{\mathrm{lin}}\)
is the longest actual linear accumulation count.

\begin{theorem}[Modular total evaluation accuracy]
\label{thm:evaluation}
Assume the half-target truncation onset, the displayed slope and
geometry gates, exact query coordinates, and exact frozen
descriptors or a certified guard-strip procedure. Suppose
certified inputs, nodes, weights, scale, and the arithmetic
recurrences satisfy
\cref{eq:allocation,eq:arithmetic-error,eq:node-data-error,eq:geometry-error,eq:RS-data-budget,eq:roundoff-updates}.
Then the separately computed field, spatial-through-three, and
first-time outputs have total outer-envelope error at most
\(\e^{-L^2/4}\).
A sufficient asymptotic requested normalized accuracy is
\begin{equation}
 P_{\mathrm{req}}=
             \frac{3L^2}{8\log2}+O_{r,b}(L\log L)
 \quad\hbox{bits}.
 \label{eq:precision}
\end{equation}
The finite displayed inequalities, rather than the \(O\)-term
alone, specify acceptance at a chosen height.
\end{theorem}
\begin{proof}
The truncation uses half the target.
The allocation proof after \eqref{eq:allocation}, followed by
\eqref{eq:outer-conversion}, bounds the evaluation error by
the other half.
For \eqref{eq:precision}, every real-node conditioning and
intermediate-size cost has logarithm \(O(L\log L)\).
Also \(\log N=L/2+O(1)\) and both quadrature orders are
\(O(L^2)\). The positive error recurrences therefore require
normalized input errors and roundoff can be made sufficiently
small by choosing them at most
\(\exp[-3L^2/8-C_{r,b}L\log L]\), with \(C_{r,b}\)
large enough for their finite majorants.
Taking base-two logarithms gives \eqref{eq:precision}.
The phase guards in \eqref{eq:phase-budget} and the small
accumulation guard are lower-order requirements within that
choice.
\end{proof}
The precision is sufficient, not necessary. It specifies
delivered input accuracy and the declared arithmetic model;
it does not bound a concrete provider's internal precision,
iteration count, or bit complexity. Those costs are supplied separately
in \cref{sec:providers,subsec:bit-complexity}.
For arbitrary real coordinates, exact floor decisions can require
arbitrarily fine discrimination near a threshold. Rounded
floor tests are not part of the theorem. Uncertainty in \(z,T\)
also needs its own derivative-based error allocation.

\subsection{Costs and the meaning of an exact residual}
With named data already supplied, polynomial preprocessing uses
\(O(p^3)\) scalar arithmetic, and all arithmetic channels use
\(O(Np)\). The weighted Riemann--Siegel table uses \(O(K^2)\).
The contour construction at all outer nodes uses
\(O(n(q_cK+K^2))=O(L^5)\) scalar arithmetic, with fixed contour
values reused. This last count is for the completion part only;
it does not replace the \(N\asymp\sqrt X\) arithmetic modes.
All counts exclude provider bit costs.

Three residual statements must remain distinct.
The exact frozen model solves heat exactly.
A symbolic model with fixed perturbed coefficients and fixed
quadrature nodes, weights, and node data also solves heat
exactly, since it is a sum of heat solutions.
Independently rounded channel outputs have only a bounded
defect: their \(T\)-minus-\(zz\) discrepancy is bounded by
the sum of the paid channel errors.
A rounded spatial derivative of that defect would require
an additional mixed \(T,z\) output channel. It is not supplied
by differentiating a pointwise rounding allowance.

\section{Finite algorithms for the evaluation data}
\label{sec:providers}
We now supply the data required by \cref{thm:evaluation}.
The construction separates delivered accuracy from the arithmetic used
to obtain it. Elementary rational series, a shifted Stirling expansion,
exact polynomial root counts, and a guarded real partition suffice.
This section integrates the finite constructions of \cite{rec09,rec10};
\cref{sec:effective} composes them with the explicit calibration below.

\subsection{Exact inputs and certified frozen descriptors}
\label{subsec:input-class}
Fix positive dyadic radii \(r,b\). A query consists of a positive dyadic
center \(X\), exact rational \(d=z-X\) and \(T\), with
\(z\in\Omega_X(r,b)\), \(1\le T\le2\), and finite rational
certificates for integers \(p,N\) satisfying
\begin{equation}
 2p<L^-\le L\le L^+<2p+2,\qquad
 N<a^-\le a_X\le a^+<N+1.
 \label{eq:descriptor-certificates}
\end{equation}
The constants are proved on the enlarged rectangle
\(\Omega_X(R,Y)\), where \(R=r+1,Y=b+1\), before obtaining the
target derivative bounds. A larger target rectangle is handled by
declaring its own radii and buffers.
Invalid certificates are rejected. Strict centers admit such witnesses
by refinement, with separation
\begin{equation}
 \sigma=\min\{L-2p,2p+2-L,a_X-N,N+1-a_X\}>0.
 \label{eq:center-separation}
\end{equation}
This is an input restriction, not an equality test for arbitrary
computable reals. No lower bound on \(\sigma\) is implicit.
All certified descriptors, coefficients, partitions, and reference
geometry are fixed under differentiation.

Let \(c_-\le c_{\mathrm o}\) and \(C_+\ge C_\chi\) be valid
positive rational calibration bounds. Choose a dyadic budget \(e>0\)
by strict enclosure tests so that
\begin{equation}
 \frac{c_-}{8C_+}\e^{-3L^2/8}<e<
 \frac{c_-}{4C_+}\e^{-3L^2/8}\le\varepsilon_X/2.
 \label{eq:comparable-budget}
\end{equation}
All positive finite majorants can be enclosed from above within a factor
of two of one plus their value. Tolerances chosen within a fixed factor
of the caps below give the stated asymptotic precision bounds. Smaller
tolerances remain valid, with their actual cost counted.

\subsection{Elementary data, logarithmic normalization, and scale}
\label{subsec:elementary-providers}
\begin{lemma}[Finite elementary enclosures]\label{lem:elementary}
Rational input admits terminating rational enclosures for \(\pi\),
positive logarithms and square roots, right-half-plane logarithms,
reduced real and complex exponentials, and reciprocals with a certified
positive separation. Every requested output tolerance can be met by
an explicit series-tail or bisection test.
\end{lemma}
\begin{proof}
The series and their remainders are given in
\cref{app:elementary}. Exact rational partial sums have no operation
rounding. Range reduction uses a freely chosen encoding integer,
then refines its multiplied input errors. The noninteger denominators
have declared positive lower bounds. Every displayed tail tends to zero.
\end{proof}
All half-integer moment constants are also effective, using
\(\Gamma(j+1/2)=(2j)!\sqrt\pi/(4^jj!)\).

For real \(u>0\), a compatible logarithm of the exact gamma factor is
\begin{equation}
 \ell_+(u)=\log\frac{1+u^2}{64}
       -\frac{1+\ii u}{4}\log\pi
       +\log\Gamma\left(\frac{1+\ii u}{4}\right)+\ii\pi.
 \label{eq:provider-logM}
\end{equation}
The other arm is its conjugate. If
\(D_m=(\log\Gamma)^{(m)}((1+\ii u)/4)/m!\), then
\begin{equation}
 \frac{\ell_+^{(m)}(u)}{m!}
 =\frac{(-1)^{m-1}}m\{(u-\ii)^{-m}+(u+\ii)^{-m}\}
   +(\ii/4)^mD_m-\mathbf1_{m=1}(\ii/4)\log\pi.
 \label{eq:provider-logM-jets}
\end{equation}
This supplies the normalized curvature coefficients and slopes without
forming derivative factorials. Use log differences to construct
\[
 \beta_{+,n}=n^{-1/2}\e^{\ii(\Ima\ell_+(X)-X\log n/2)},\qquad
 \frac{M(w_+(u))}{g}=\e^{\ell_+(u)-\Rea\ell_+(X)}.
\]
The phase budget \eqref{eq:phase-budget} remains in force; in particular
the error in \(\log n\) is multiplied by \(X/2\).

For the common scale write \(g=\e^t\), \(t=\Rea\ell_+(X)\).
Choose an integer \(k\) from a coarse enclosure of \(t/\log2\),
then certify \(\abs{t-k\log2}\le2\). Deliver both the log-scale
error and the relative reduced-exponential error below
\(\varepsilon_g/8\). For \(\varepsilon_g\le1/2\) their product
gives \(\abs{\widehat g/g-1}\le\varepsilon_g\).
Return \(\widehat g=\widehat m2^k\), with unrestricted integer exponent.
Here \(\abs k=O(X+\log X)\), but its binary length is \(O(\log X)\).
The accuracy requested for \(\log2\) includes its multiplier \(\abs k\).

\subsection{A short gamma construction with simultaneous jet errors}
\label{subsec:short-gamma}
\begin{proposition}[Shifted Stirling provider]\label{prop:gamma-provider}
Let \(u>0\) be rational with binary input length \(b_\Gamma\).
For integers \(P\ge1,J\ge0\), normalized log-gamma derivatives
through order \(J\), including the value, can be delivered with
error at most \(2^{-P}\). With
\(t_\Gamma=P+J+b_\Gamma+10\), sufficient bounds are
\[
 O(t_\Gamma^6)\quad\text{bit operations},\qquad
 O(t_\Gamma^3)\quad\text{bits of storage}.
\]
The expansion and recurrence use \(O(P+J)\) terms and shifts.
\end{proposition}
\begin{proof}
Put \(\zeta=(1+\ii u)/4\), \(m=P+2J+4\), \(a=m+1\),
and \(W_\Gamma=\zeta+a\). Define
\begin{equation}
 S_m(w)=(w-\tfrac12)\Log w-w+\tfrac12\log(2\pi)
       +\sum_{k=1}^{m-1}\frac{B_{2k}}{2k(2k-1)w^{2k-1}}.
 \label{eq:short-stirling}
\end{equation}
The complex log-gamma remainder in \cite[\S5.11(ii)]{dlmfgamma}
is bounded by the first omitted term times
\(\sec^{2m}(\arg w/2)\). The Bernoulli Fourier formula
\cite{dlmfbernoulli} gives \(\abs{B_{2m}}\le4(2m)!/6^{2m}\).
For \(\Rea w\ge m\), consequently,
\begin{equation}
 \abs{\log\Gamma(w)-S_m(w)}
 \le \frac{2^m4(2m)!}{6^{2m}2m(2m-1)m^{2m-1}}
 \le\frac{(2/9)^m}{m}\le\frac{4^{-m}}m.
 \label{eq:short-gamma-remainder}
\end{equation}
This bound holds throughout \(w=W_\Gamma+s\), \(\abs s\le1/8\).
The analytic gamma recurrence, with logarithms fixed from positive real
values, gives
\[
 \log\Gamma(\zeta+s)=S_m(W_\Gamma+s)
       -\sum_{h=0}^{a-1}\Log(\zeta+h+s)+R_m(W_\Gamma+s).
\]
There is no principal-log phase reset. Cauchy's formula gives, for
\(0\le j\le J\),
\begin{equation}
 \abs{[s^j]R_m(W_\Gamma+s)}
       \le8^j4^{-m}/m\le2^{-P-8}.
 \label{eq:gamma-jet-remainder}
\end{equation}
The remaining coefficients are finite rational expressions and logarithms:
\[
 [s^j](w+s)^{-(2k-1)}
   =(-1)^j\binom{2k+j-2}{j}w^{-2k-j+1},\quad
 [s^j]\Log(w+s)=\frac{(-1)^{j-1}}{jw^j}\quad(j\ge1).
\]
Writing \(L_0=\Log W_\Gamma\),
\(L_j=(-1)^{j-1}/(jW_\Gamma^j)\), the leading Stirling part has
coefficients \((W_\Gamma-1/2)L_0-W_\Gamma+\log(2\pi)/2\),
\((W_\Gamma-1/2)L_1+L_0-1\), and
\((W_\Gamma-1/2)L_j+L_{j-1}\) for \(j\ge2\).

Here is one complete internal allowance. Put \(\epsilon=2^{-P}\),
\(T_\Gamma=a+m+J+10\), and
\(U_\Gamma=a+1+(\abs u+1)/4\ge\abs{W_\Gamma}\).
Deliver \(\Log W_\Gamma\), \(\log(2\pi)\), and each shift logarithm
to errors \(\epsilon/[64(1+U_\Gamma)]\), \(\epsilon/32\), and
\(\epsilon/(64a)\), respectively. Compute each rational term exactly,
quantize its complex value to error \(\epsilon/(64T_\Gamma)\), and
sum on a common dyadic grid. Fewer than \(T_\Gamma\) terms enter any
jet. These errors, \eqref{eq:gamma-jet-remainder}, and final quantization
sum to less than \(\epsilon\).

All shift denominators have modulus at least \(1/4\). The binomial
factor divided by \(\abs{W_\Gamma}^j\) is at most \(3^j/j!<27\);
the correction terms are bounded by \(27/(m9^k)\), the shift jets
by \(a4^J\). A sufficient internal dyadic width is
\(P+O(\log(2+\abs u)+J+\log(P+J+2))\).
Bernoullis are generated by their triangular rational recurrence;
their denominators divide \((n+1)!\). Cached powers and elementary-series
denominators have \(O(t_\Gamma^2)\) bits. There are
\(O(t_\Gamma^2)\) integer operations, each costing
\(O(t_\Gamma^4)\) with elementary multiplication and division.
Processing shift arrays sequentially gives the stated cubic space.
Exact temporary operands can exceed the delivered precision.
\end{proof}
The elementary Euler-product alternative in \cref{app:euler-provider}
also proves finite delivery, but needs a much longer sum.
Its accumulation error explains why output accuracy alone cannot be
used as a provider's internal roundoff allowance.

\subsection{Exact Gauss root counts and polynomial work}
\label{subsec:gauss-provider}
Let \(P_s\) be the Rodrigues Legendre polynomial. Orthogonality implies
\(s\) simple roots in \((-1,1)\): otherwise the product of its fewer
sign-changing factors would contradict orthogonality.
Define \(V_s(x)\) as the sign variation of
\((P_s(x),\ldots,P_0(x))\), with zeros omitted.
The identities
\[
 (k+1)P_{k+1}=(2k+1)xP_k-kP_{k-1},\qquad
 (1-x^2)P_s'=s(P_{s-1}-xP_s)
\]
follow from Rodrigues' formula; see also \cite{dlmflegendre}.
At an internal-chain zero the neighbors have opposite signs, so variation
is unchanged. At a root of \(P_s\) it drops by one, and omitting the
leading zero gives its right limit. Hence the exact root count in an
open interval is
\begin{equation}
 \#\{\xi_i\in(A,B)\}
 =V_s(A)-V_s(B)-\mathbf1_{P_s(B)=0}.
 \label{eq:legendre-open-count}
\end{equation}
Bisect positive-count dyadic intervals, record midpoint roots as
singletons, and discard empty children. At most \(s\) intervals remain
at each depth, giving \(O(sh)\) count queries through depth \(h\).
For \(x=c/q\), compute signs using integers
\begin{equation}
 b_0=1,\quad b_1=2c,\quad
 b_{k+1}=2(2k+1)c b_k-4k^2q^2b_{k-1},
 \quad b_k=q^k2^kk!P_k(c/q).
 \label{eq:integer-legendre}
\end{equation}
Thus no rational sign oracle or unhandled endpoint equality is needed.

Define the polynomial
\(J_s(a)=\int_{-1}^1(P_s(t)-P_s(a))/(t-a)\dd t\).
The exact Gauss weights are
\begin{equation}
 \omega_i=J_s(\xi_i)/P_s'(\xi_i).
 \label{eq:rational-gauss-weight}
\end{equation}
Polynomial division proves exactness through degree \(2s-1\), while
the nodal Lagrange polynomial gives
\(\omega_i=\int_{-1}^1l_i^2>0\) and \(\sum_i\omega_i=2\).
Explicit separation and conditioning constants are recorded in
\cref{app:gauss-conditioning}; they prove the following bound.

\begin{proposition}[Finite Gauss data]\label{prop:gauss-provider}
For base-node error \(\eta_x\), total base-weight error \(\eta_w\),
and rational input length \(\ell_{\rm in}\), put
\[
 U_s=P_{\rm samp}+\ell_{\rm in}
                 +s^3\lceil\log_2(s+2)\rceil+10,
\]
where \(P_{\rm samp}\ge1\) bounds both negative binary log tolerances.
Positive rational weights and interior rational nodes can be delivered
with those errors in \(O(s^4U_s^3)\) bit work and \(O(s^2U_s)\) space.
The required isolation depth is
\(O(P_{\rm samp}+\ell_{\rm in}+s^3\log(s+2))\).
\end{proposition}
\begin{proof}
Use \eqref{eq:legendre-open-count} until the explicit root/weight
allowances in \cref{app:gauss-conditioning} hold. Each variation query
uses \(O(s)\) operations on \(O(s(h+\log(s+2)))\)-bit integers.
Exact homogeneous weight evaluation and coefficient construction fit
the same work and space bounds. A singleton root is handled exactly.
\end{proof}

\subsection{Guarded geometry at cutoff transitions}
\label{subsec:guard-geometry}
Write \(\alpha=X-V,\beta=X+V,\vartheta_j=4\pi j^2\).
For rational \(0<\eta\le\min(\varepsilon_b,1/4)\), compute strict
rational enclosures of width at most \(\eta\) for \(\alpha,\beta\)
and \(\vartheta_N,\vartheta_{N+1}\). Padding supplies strictness even
at a rational endpoint. A threshold is a candidate if its enclosure
intersects the enclosing window. The local-strip gates imply
\[
 a'(u)\le\frac1{2\sqrt{2\pi X}},\qquad
 \operatorname{span}_{[X-(V+1),X+(V+1)]}a
 \le\frac{V+1}{\sqrt{2\pi X}}\le\frac18.
\]
There is therefore at most one candidate, and every true threshold
is a candidate. Only \(N,N+1\) need be considered.
With enclosures \([\alpha^-,\alpha^+]\), \([\beta^-,\beta^+]\),
\([\vartheta_j^-,\vartheta_j^+]\), retain
\begin{equation}
 [\alpha^+,\min(\beta^-,\vartheta_j^-)],\quad
 [\max(\alpha^+,\vartheta_j^+),\beta^-]
 \label{eq:interior-branches}
\end{equation}
with labels \(j-1,j\). Without a candidate retain
\([\alpha^+,\beta^-]\) with label \(N\).
Discard nonpositive intervals and pieces of length at most \(\eta\).
Endpoint and threshold trimming omit at most \(3\eta\); discarding
short pieces adds at most \(2\eta\). Thus the total missing integral
is bounded by \(5\eta B_{\rm loc}\le6\varepsilon_b B_{\rm loc}\).
Every retained piece is inside the correct original branch, including
when a threshold equals a window edge or lies just outside it.

Apply the Gauss rule by the exact rational affine map
\(u=c+h\xi\), \(\mu=h\omega\), with \(0<h\le V\).
Base errors give physical errors at most \(V\eta_x,V\eta_w\);
split the total weight allowance across the at most two branches.
Each subinterval ellipse lies inside its parent ellipse, so
\eqref{eq:quadrature-bound} still applies. No division by branch length
occurs, no rounded floor relabels a node, and trimming is charged once.

\subsection{Pole-safe coefficient delivery}
\label{subsec:pole-safe-delivery}
Use the contour formula \eqref{eq:contour-jet}, the exact weighted
recurrence \eqref{eq:weighted-rs}--\eqref{eq:rs-endpoint}, and all
\(m\le M:=3K\). The contour has length \(12\), denominator
\(\abs{2\cos(\pi\zeta)}\ge1\), and target separation at least
\(1/4\). The entire continuation, maximum modulus, and Cauchy give
\[
 B=B_c4^{M+1},\qquad L_\zeta=4(M+2)B,\qquad L_\chi=4(M+1)B
\]
as common magnitude and Lipschitz bounds for its integrands.
Keep \(\widehat\chi\) within \(1/8\) of the exact real value.
If the contour's total weight, node, target-coordinate, primitive, and
\(\kappa=1/(2\pi\ii)\) errors are
\(\eta_w,\eta_\zeta,\eta_\chi,\eta_H,\eta_\kappa\), the additional
error in any normalized jet is at most
\begin{equation}
 \frac{\eta_wB+(12+\eta_w)
       (L_\zeta\eta_\zeta+L_\chi\eta_\chi+\eta_H)}{2\pi}
 +\eta_\kappa(12+\eta_w)(B+\eta_H).
 \label{eq:contour-delivery}
\end{equation}
For common errors at most \(\delta\le1/8\), bound this by
\(\delta C_c\), where
\(C_c=[B+13(L_\zeta+L_\chi+1)]/6+13(B+1)\).
Choose \eqref{eq:contour-error} below \(\tau_h/4\) and
\(\delta\le\min(1/8,\tau_h/[4(1+C_c)])\); final quantization can
use another \(\tau_h/4\). At contour nodes, numerator and denominator
errors give \(\abs{\widehat F_0-F_0}\le2\eta_N+2B_c\eta_D\)
when \(\eta_D\le1/2\). Requiring this below
\(\eta_H4^{-M-1}\) supplies the primitive allowance.

For explicit coefficient propagation put
\begin{align*}
 A_{kj}&=\abs{e_j^{(k)}}2^{-j}3^{j-2k},&
 S_C&=\max_{k\le K}\sum_jA_{kj},\\
 T_C&=\max_{k\le K}\sum_j\frac{2k-j}{3}A_{kj}
                         (B_F2^{3k-2j}+1).
\end{align*}
Since \(j\le3k/2\), the powers of \(\pi\) have nonpositive
exponents. The mean value theorem, \(\pi,\widehat\pi\ge3\), and
\(\tau_h\le1\) imply
\begin{equation}
 \abs{\widehat C_k-C_k}\le\tau_hS_C+
                           \abs{\widehat\pi-\pi}T_C.
 \label{eq:Ck-delivery}
\end{equation}
Taking the two terms below \(\varepsilon_C/4\) leaves quantization
slack. Both condition sums have logarithms \(O(K\log(K+2))\).
This constructs every coefficient, including at removable poles.
It uses the named exact identities, with their source-review scope
stated in \cref{sec:quadrature}; no preliminary precision assertion
from the coefficient literature is imported.

\section{Explicit calibration and an all-height onset}
\label{sec:calibration}
The constants used in the preceding comparisons can be chosen
explicitly. The point of this section is their uniform domains:
a scalar evaluation at a large height would not prove a disk bound,
an all-order coefficient estimate, or a statement through every change
of \(p\) and \(N\). We give the calibration of \cite{rec11} and its
continuous majorants before using it in the effective construction.

Set
\[
 W=R+Y,\qquad d_{\rm ch}=\lceil W\rceil,\qquad
 \nu=d_{\rm ch}+1\ge4.
\]
The chart integer \(d_{\rm ch}\) is distinct from the displacement
\(d=z-X\). A sufficient rational record is
\begin{equation}
\begin{gathered}
 H=A=C_\Gamma=c_{\mathrm o}=1,\quad D=4,\quad D_0=3,
 \quad g_-=1/64,\\
 C_H=M_H^*=1/4,\quad C_{\rm tail}=640,\quad
 C_\chi=3^{d_{\rm ch}},\quad C_{\rm cal}=32+d_{\rm ch},
 \quad X_{\rm cl}=64\nu.
\end{gathered}
\label{eq:explicit-calibration}
\end{equation}
Here \(M_H^*\) bounds the complete source mass \(M_H\), rather than
denoting its exact value. The domain statements follow.

\subsection{A complex-disk estimate before differentiation}
\label{subsec:disk-calibration}
The digamma remainder \cite[\S5.11(ii)]{dlmfgamma} gives
\begin{equation}
 \psi(t)=\Log t-\frac1{2t}+E(t),\qquad
 \abs{E(t)}\le\frac{\sec^3(\arg t/2)}{12\abs t^2},
 \quad\abs{\arg t}<\pi.
 \label{eq:digamma-calibration}
\end{equation}
For \(t_\pm=(1\pm\ii z)/4\) on \(\abs{z-X}\le X/2\),
one has \(\abs t\ge X/8\), \(\abs{\arg t}<3\pi/4\).
Since \(\sec(3\pi/8)<3\), \(\abs E\le144/X^2\).
Each point of the inner \(X/4\) disk has a \(t\)-disk of radius
\(X/16\) in this larger region. Thus \(\abs{E'}\le2304/X^3\).
Using \(\abs t\ge3X/16\) and
\(\abs w,\abs{w-1}\ge3X/8\) on the inner disk in the exact
formula for \(\ell_\pm''\) proves
\begin{equation}
 \sup_{\abs{z-X}\le X/4}\abs{\ell_\pm''(z)}
 \le\frac1{3X}+\frac{40}{9X^2}+\frac{144}{X^3}
 \le\frac1X\qquad(X\ge32).
 \label{eq:explicit-disk}
\end{equation}
At \(X=32\), multiplication by \(X\) gives \(353/576<1\);
the bound decreases thereafter. Cauchy's estimate now supplies all
orders simultaneously:
\[
 \abs{c_{\epsilon k}}\le
       \frac{4^{k-1}}{k(k+1)}X^{-k}\le4^{k-1}X^{-k}.
\]

At a real center, direct differentiation of
\eqref{eq:provider-logM} yields
\begin{align*}
 \ell_+'(X)&=-\pi/8+\ii L/4+\delta(X),\\
 \delta(X)&=\frac{2X}{1+X^2}+\frac\ii4\Log(1-\ii/X)
             -\frac\ii{8t}+\frac\ii4E(t),\quad t=(1+\ii X)/4.
\end{align*}
Here \(\abs E\le4/X^2\), and for \(X\ge32\),
\begin{equation}
 \abs\delta\le\frac2X+\frac1{4(X-1)}+
                  \frac1{2X}+\frac1{X^2}<\frac3X.
 \label{eq:explicit-slope}
\end{equation}
Reflection gives the other arm. Along real \(u\ge32\),
\(\abs{(\log G)'(u)}\le\pi/8+3/u<1\). Hence
\(G(X+q)\le g\e^{\abs q}\) for real \(\abs q\le X/2\),
\(X\ge64\), proving \(C_\Gamma=1\).

Stirling's formula with its first omitted-term bound gives exactly
\begin{equation}
 G(X)=\frac{\pi^{1/4}}{32}(1+X^2)^{7/8}\e^{-\pi X/8}
  \exp\left\{\frac X4\arctan(1/X)-\frac14+\Rea R_\Gamma(t)\right\},
 \quad\abs{R_\Gamma(t)}\le\frac2{3X}.
 \label{eq:exact-real-gamma-calibration}
\end{equation}
The final exponent exceeds \(-1/2\) for \(X\ge32\), giving
\(g\ge X^{7/4}\e^{-\pi X/8}/64\).
For \(X\ge\max(64,4W)\), integrate \((\log G)'\) from the real
center through the buffered rectangle and use \eqref{eq:explicit-disk}.
Its modulus stays below one, so both
\(g/\abs{G(z)}\) and \(\abs{G(z)}/g\) are at most
\(\e^W\le3^{d_{\rm ch}}\).

\subsection{Bounds for the complete real source}
The beta integral and half-gamma modulus \cite{dlmfbeta,dlmfhalf}
give, for real \(v\),
\[
 \left|\frac{\Gamma(1/4+\ii v)}{\Gamma(1/2+\ii v)}\right|
 \le\frac{\Gamma(1/4)}{\Gamma(1/2)},\qquad
 \abs{\Gamma(1/4+\ii v)}<8\e^{-\pi\abs v/2}.
\]
For the last inequality split the positive integral for \(\Gamma(1/4)\)
at one to obtain \(\Gamma(1/4)<5\), and use
\(\cosh(\pi v)\ge\e^{\pi\abs v}/2\).
The finite-sum zeta identity \cite{dlmfzeta} at \(N=1\) reads
\[
 \zeta(s)=\frac{s}{s-1}-s\int_1^\infty\{x\}x^{-s-1}\dd x.
\]
For \(s=(1+\ii u)/2\), it gives
\(\abs{\zeta(s)}\le1+2\abs s\le2(1+\abs u)\).
The exact normalization \(H(0,u)=M(s)\zeta(s)\) therefore yields
\begin{equation}
 \abs{H(0,u)}\le\tfrac14(1+\abs u)^4\e^{-\pi\abs u/8}
 \qquad(u\in\R).
 \label{eq:global-source-calibration}
\end{equation}
Every summand of \(\Phi(v)\) is positive for \(v\ge0\), since
\(2\pi n^2\e^{4v}>3\). Thus, without a source cutoff,
\begin{equation}
 M_H=\int_0^\infty\abs\Phi=H(0,0)\le1/4.
 \label{eq:complete-source-mass}
\end{equation}
For \(\abs q\le X/2\), \(X\ge32\),
\(1+X+q\le(5/2)X\). Dividing
\eqref{eq:global-source-calibration} by the lower gamma scale gives
\[
 \frac{\abs{H(0,X+q)}}g
     \le640X^{9/4}\e^{-\pi q/8}.
\]
Consequently the moderate and remote constants can be taken as
\begin{equation}
 A_{\rm mod}^{\theta}=640\max(\mathsf D_0^\theta,
                    (\mathsf D_0^\theta)^t),\qquad
 A_{\rm far}=16\max(\mathsf J_0,\mathsf J_0^t).
 \label{eq:explicit-tail-calibration}
\end{equation}

\subsection{Calibration of the actual outer envelope}
Put \(L_z=\Log(z/(4\pi))\),
\(G_0(z)=\pi^{1/4}z^{7/4}\e^{-\pi z/8}/32\), and
\(\theta_0=z(L_z-1)/4+7\pi/8\).
On the buffered rectangle, when \(X\ge\max(16,2(W+2))\),
the exact sectorial formula gives
\begin{equation}
 \ell_\pm(z)=\log G_0(z)\pm\ii\theta_0(z)+\delta_\pm(z),
 \qquad\abs{\delta_\pm(z)}\le8/\abs z\le16/X.
 \label{eq:complex-gamma-calibration}
\end{equation}
For the plus arm the remainder is
\[
 \Log(1+z^{-2})+\frac{\ii z-1}{4}\Log(1-\ii/z)
                      -\frac14+R_\Gamma(t_+).
\]
The logarithm series and \(\abs{R_\Gamma(t_+)}\le16/(3\abs z)\)
bound this by \(8/\abs z\); conjugate reflection gives the minus arm.
The branches are fixed at positive real height. Averaging preserves the
canonical \(G\) and gives \(\abs{\log(G/G_0)}\le16/X\).

The original amplitude is
\(A_{\mathrm o}(-T,z)=2G_0(z)\e^{-TL_z^2/16+\pi^2T/64}\).
Since \(\abs{L_z-L}\le2W/X\),
\(\abs{L_z^2-L^2}\le4W(1+L)/X\) for \(L\ge0\), one obtains
\begin{equation}
 \left|\log\frac{\abs{A_{\mathrm o}(-T,z)}}
 {2\abs{G(z)}\e^{-TL^2/16+\pi^2T/64}}\right|
 \le\frac{C_{\rm cal}}2\frac{1+L}{X}.
 \label{eq:explicit-envelope-calibration}
\end{equation}
Under \(X\ge X_{\rm cl}\) and
\(C_{\rm cal}(1+L)/X\le1\), AM--GM for the two original arms gives
\begin{equation}
 \calE_{\mathrm o}(-T,z)\ge\abs{A_{\mathrm o}(-T,z)}
 \ge\abs{G(z)}\e^{-TL^2/16+\pi^2T/64}.
 \label{eq:explicit-envelope-lower}
\end{equation}
The amplitude ratio in \eqref{eq:explicit-envelope-calibration}
differs from one by at most \(C_{\rm cal}(1+L)/X\).
This proves \(c_{\mathrm o}=1\) while retaining the transverse growth
of the actual two-arm envelope. The common center onset \(64\nu\)
covers the full gamma disk, real window, and buffered chart;
the remaining height-dependent structural gates are paid next.

\subsection{Continuous majorants through every order jump}
\label{subsec:explicit-onset}
For clarity first retain a valid general calibration record.
Let \(\theta=1/2\), \(C_{\mathrm o}=C_\chi/c_{\mathrm o}\),
and choose fixed constants \(S,M\) so that
\begin{align}
 \log\mathcal J(p),\log\mathcal J_\theta(p)
   &\le S(p+1)\log(p+2),\qquad S\ge1,\nonumber\\
 M&=\max\{0,\log^+(8C_{\mathrm o}C_Q),
       \log^+(392C_{\mathrm o}C_{\rm pop}),\nonumber\\
 &\hspace{36mm}\log^+(392C_{\mathrm o}A_{\rm mod}^\theta),
       \log^+(392C_{\mathrm o}A_{\rm far})\},
 \label{eq:onset-log-costs}
\end{align}
where \(Q_p\le C_Q\mathcal J(p)\) and
\(A_{\rm art,p}^\theta\le C_{\rm pop}\mathcal J_\theta(p)\).
For example \(C_Q=48M_0+M_0^t+49C_{\rm pop}\e\) works.
The factorial estimates in \cref{sec:growing} supply these constants.
With \(A=\max(1,H),D=4\), a sufficient general onset is
\begin{align}
 L_{\rm rate}&=\max\{16,(32S)^2,\sqrt{32M}\},\nonumber\\
 L_0&=\max\bigl\{L_{\rm rate},(2R+W+1)/8,\log(1+4(W+1)),
       2\log(1+24A),\nonumber\\
 &\hspace{22mm}2\log(1+D_0),2\log(1+C_{\rm cal}),
                  \log(1+X_{\rm cl}/(4\pi))\bigr\}.
 \label{eq:parameterized-onset}
\end{align}
Here \(X_{\rm cl}\) must already be a valid \emph{center} onset
for every classical domain, including the real window.

Indeed, for \(p=\lfloor L/2\rfloor\), \(L\ge16\),
\[
 (p+1)\log(p+2)\le L\log(L+2)\le L^{3/2},\qquad
 (p+3/4)\log X-3L^2/8\ge7L^2/64.
\]
At \(L\ge L_{\rm rate}\), \(M+SL^{3/2}\le L^2/16\),
so the local half-target gate holds strictly.
The moderate and artificial gates have available quadratic coefficient
\(61/8\); their height powers and logarithmic costs fit with slack.
For the remote gate, \(\pi X/8\le X^2/512\) and
\(X^2/512\ge9L^2/8>(3/8+1/32)L^2\).
These are the four conditions \eqref{eq:growing-gates} with factors
\(8,392\), reserving half the total target for evaluation.

The structural conditions also follow from \eqref{eq:parameterized-onset}.
For \(L\ge16\), \(V\le12L\), \(X\ge12\e^L\), and
\(\sqrt{2\pi X}\ge8\e^{L/2}\). Hence
\[
 \frac{4(V+2)}{\sqrt{2\pi X}}\le(6L+1)\e^{-L/2}<\frac12,
 \quad2(V+2)\le X/4,\quad K\le L\le(X-V)/(4\pi).
\]
The buffer terms in \(L_0\) give the window gates;
\(L\ge2\log(1+24A)\) pays \(2AV^2/X\le1\).
The slope terms give \(X\ge10D_0\) and \(X\ge2D_0(L+2)\),
which pay both the stronger real-slope gate and the separate damping
condition in \eqref{eq:structural-gates}.
The calibration term pays \(C_{\rm cal}(1+L)/X\le1\).
Strong geometry gives the remaining ellipse/logarithm bounds.
No monotonicity of a factorial expression across floor jumps is assumed.

\begin{theorem}[Explicit sufficient onset]\label{thm:explicit-onset}
For the record \eqref{eq:explicit-calibration}, let
\begin{equation}
 L_*=2^{22}\nu^4,\qquad X_*=2^{2L_*+8}.
 \label{eq:closed-onset}
\end{equation}
Every center with \(L\ge L_*\) meets the uniform classical domains,
structural gates, and half-target truncation gates for
\(p=\lfloor L/2\rfloor,K=2p,1\le T\le2\).
The condition \(X\ge X_*\) is a stronger explicit dyadic sufficient
condition. Finite evaluation additionally retains
\eqref{eq:descriptor-certificates} and its separation dependence.
\end{theorem}
\begin{proof}
Use the rational majorants
\[
 \overline E=3^{8\nu^2},\quad\overline S=32\nu,\quad
 \overline Q=32\overline S^2=2^{15}\nu^2,\quad
 \overline T=\overline E\,\overline S^2,
\]
and \(\overline{\mathcal J}(p)=
\overline T(p+2)\overline Q^{p+1}(2p+4)!\).
They dominate both moment systems in \cref{sec:growing}:
\(1+R/4\le\nu/2\), the largest refined-tail exponent is at most
\(9\nu^2/4\), and \(\overline S\) dominates every moment base.
Also the RS constants satisfy \(Q_R<64,C_R<32\).
The factorial bound gives
\[
 \log\overline{\mathcal J}(p)
 \le(6+\log\overline Q+\log\overline T)(p+1)\log(p+2),
 \quad\log\overline T\le18\nu^2,\quad\log\overline Q\le2\nu^2.
\]
The explicit moments imply
\(M_0\le\overline E\), \(M_0^t\le\overline T\),
\(\mathsf D_0^\theta,(\mathsf D_0^\theta)^t,
\mathsf J_0,\mathsf J_0^t\le\overline T\).
For the remote moments one may use
\(\mathsf J_m\le2\e^{Y^2/4}m!(Y+4)^m\).
It follows that
\[
 C_Q\le48\overline E+\overline T+588\le3\overline T,
 \quad A_{\rm mod}^\theta\le640\overline T,
 \quad A_{\rm far}\le16\overline T.
\]
Since \(C_{\mathrm o}=3^{d_{\rm ch}}\), every fixed logarithmic
cost in \eqref{eq:onset-log-costs} is bounded using a multiplier
below \(2^{18}\). Thus \(S,M\le32\nu^2\), and
\(L_{\rm rate}\le2^{20}\nu^4\).
All remaining terms of \eqref{eq:parameterized-onset} are smaller than
\(2^{22}\nu^4\). Finally \(\log2>1/2\), \(\log(4\pi)<3\)
give \(\log(X_*/(4\pi))>L_*+1\).
\end{proof}
The threshold is deliberately conservative. Its role is effectiveness,
not a practical starting height. The radii and any inherited regulator
dependence are fixed before height grows. Neither an endpoint-uniform
time theorem nor a uniform lower bound on center separation follows.

At a smaller specified height there is also a finite verifier: recompute
outward rational enclosures for \eqref{eq:descriptor-certificates}, all
structural/domain gates, and the four half-target inequalities, and
accept only certified strict margins. A finite unsuccessful witness is
rejected or inconclusive. Every strictly admissible height has a finite
successful witness. This does not promote a single-height verification
to an all-height theorem.

\subsection{A smaller onset from additive costs}
\label{subsec:linear-onset}
The preceding onset is a valid fallback. Its main loss comes from
multiplying the fixed chart cost \(\log\overline T\) by
\((p+1)\log(p+2)\), followed by \(L\log(L+2)\le L^{3/2}\).
Keeping that cost additive gives the following improvement
\cite{rec16}. The moment and factorial majorants themselves need not change.

\begin{theorem}[Linear sufficient log-height]\label{thm:linear-onset}
With the same calibration record and frozen choices, every center with
\begin{equation}
 L\ge L_{\rm new}=64\nu
 \label{eq:linear-onset}
\end{equation}
meets all classical, structural, calibration and half-target truncation
gates of \cref{thm:explicit-onset}. A stronger dyadic sufficient condition is
\begin{equation}
 X\ge X_{\rm new}=2^{128\nu+8}.
 \label{eq:linear-dyadic-onset}
\end{equation}
All input and descriptor restrictions remain. In particular the theorem
does not admit excluded boundary centers or bound their separation uniformly.
\end{theorem}
\begin{proof}
Keep \(\overline E,\overline S,\overline Q,\overline T\) and
\(\overline{\mathcal J}(p)\) from the preceding proof.
They dominate both moment systems, with
\[
 Q_p\le C_Q\overline{\mathcal J}(p),\quad C_Q\le3\overline T,
 \quad A_{\rm art,p}^{1/2}\le C_{\rm pop}\overline{\mathcal J}(p),
 \quad \log\overline T\le18\nu^2,\quad
 \log\overline Q\le15+2\nu\le23\nu/4.
\]
The superscript \(1/2\) labels \(\theta\).
For every \(L\ge64\nu\ge256\),
\begin{equation}
 \log(L+4)\le L/32,\quad p+1\le129L/256,\quad
 2p+4\le65L/64.
 \label{eq:linear-elementary-bounds}
\end{equation}
Indeed \(\e^8>260\) by the exponential series through degree four,
and \(L/32-\log(L+4)\) increases on this interval.
Since \(\log24+(\nu-1)\log3\le\nu^2\),
the local logarithmic cost satisfies
\begin{align}
 \Lp(8C_{\mathrm o}Q_p)
 &\le37\nu^2+\log(p+2)+(p+1)\log\overline Q
                         +\log((2p+4)!)\nonumber\\
 &\le\frac{592+8+2967+2080}{65536}L^2
   <\frac3{32}L^2.
 \label{eq:linear-local-cost}
\end{align}
Here \(m!\le m^m\) pays the factorial, and the four rational
summands bound the four terms in the first line, respectively.
Also \(8C_Q\ge392C_{\rm pop}\e>392C_{\rm pop}\), so the same
bound applies to \(\Lp(392C_{\mathrm o}A_{\rm art,p}^{1/2})\).
The unchanged floor lower bound gives
\[
 (p+3/4)\log X-3L^2/8\ge7L^2/64>3L^2/32,
\]
paying the local half-target gate with strict slack.

For the true tails, the preceding moment estimates give
\[
 \Lp(392C_{\mathrm o}A_{\rm mod}^{1/2}),\quad
 \Lp(392C_{\mathrm o}A_{\rm far})
 \le18+2\nu+18\nu^2\le20\nu^2<L^2/32.
\]
Since \(\log X<L+3\), the moderate height power costs at most
\(L^2/4\), and the artificial one at most \(L^2/16\).
Their combined right sides are below \(9L^2/32\) and \(5L^2/32\),
respectively, both smaller than \(61L^2/8\).
Finally \(X>24L\), \(X>64\pi\), whence
\[
 \pi X/8<X^2/512,\qquad
 X^2/512>9L^2/8>(3/8+1/32)L^2.
\]
This pays the remote half-target gate. The factors \(8,392\),
and hence the half-target reservation for evaluation, are unchanged.

For completeness every non-rate gate holds on this smaller domain.
Use \(W\le\nu-1<L/64\), \(V\le12L\) and
\(X>12\e^L\ge2L^3\).
These give \(X>X_{\rm cl}=64\nu\), \(X>4(W+1)\),
\(V<X/4\), and \(2(V+2)<X/4\).
The strong ellipse and cutoff-span estimates already proved for
\(L\ge16\) apply. Also
\[
 K\le L\le(X-V)/(4\pi),\quad
 DV/X<24/L^2<1/2,\quad 2AV^2/X<128/L\le1/2,
\]
and \(V>\max(2R,W+1,64)\).
Since \(X>30\) and \(\pi<16/5\), the stronger slope gate holds.
The separate damping gate follows from
\[
 2\{(3/X)(\pi/4+L/2)+9/X^2\}
 <3(L+2)/X+18/X^2<1,
\]
using \(X>6(L+2)\) and \(X>6\).
Finally \(C_{\rm cal}=\nu+31<L\), so
\(C_{\rm cal}(1+L)<2L^2<X\).
Strong geometry supplies the ellipse logarithm bound as before.
Provider tolerances and strict descriptor witnesses retain their
separate finite contracts.

All estimates hold continuously for every \(L\ge64\nu\), using
\(L/2-1\le p\le L/2\), and the full-population bounds and exact cutoff
transfer cover every later \(N\) jump. No monotonicity of actual
factorial costs, or differentiation across chart changes, is asserted.
Lastly \(\log2>1/2\), \(\log(4\pi)<3\) give
\(\log(X_{\rm new}/(4\pi))>64\nu+1\).
\end{proof}

The old and new sufficient log-heights have exact ratio
\(L_*/L_{\rm new}=2^{16}\nu^3\ge2^{22}\).
For example on the dyadic chart family \(r=b=(\nu-3)/2\),
\(\nu=4,8,16\) gives new bounds \(256,512,1024\), against
\(2^{30},2^{34},2^{38}\).
These are symbolic comparisons of sufficient bounds, not field
measurements. The proof retains every arithmetic mode and completion
coefficient, both arms, the bounded integral, all buffers and channels,
and the fixed time interval. \Cref{sec:moving-endpoint} separately proves
\([1/L,2]\) at sufficient \(1024\nu^2\), with its own kernel and delivery
payments. Optimal onset, useful runtime, strong precision through zero
and complete evaluator certification remain separate questions.

\section{Effective evaluation of the canonical field}
\label{sec:effective}
We combine the approximation, finite data constructions, and explicit
calibration into a single algorithm theorem. The exact-input restriction,
the full arithmetic population, and the cost of small separation remain
part of the statement. The numerical implementation status is separate. The expanded-domain
construction in \cref{sec:moving-endpoint} retains this section as the
\([1,2]\), \(64\nu\) fallback; its altered outer-node and witness costs are
stated in \cref{subsec:moving-cost}.

\subsection{A finite allocation that meets every contract}
\label{subsec:finite-allocation}
Use \eqref{eq:explicit-calibration}, the comparable budget
\eqref{eq:comparable-budget}, and rational upper enclosures of the
positive quantities below. A bar denotes an upper enclosure; it is
suppressed when this creates no ambiguity. Each can be chosen below
twice one plus the exact quantity.

\begin{enumerate}[leftmargin=*,itemsep=.5em]
\item \textbf{Common scale.}
Take \(\varepsilon_g\le\min(1/2,e/[8(1+A_0)])\).
Then \(\varepsilon_gA_0\le e/8\).

\item \textbf{Arithmetic inputs.}
Put
\[
 C_{\rm ar}=C_{\rm pop}X^{1/4}
       [\mathcal D_p+2(\mathcal D_p^{(1)}+\mathcal C_p+\mathcal S_p)],
 \qquad\delta_{\rm ar}\le\min(1/4,e/[80(1+C_{\rm ar})]).
\]
Deliver \(\beta\) with error \(\delta_{\rm ar}n^{-1/2}\), and
\(\lambda,c_r\) with error \(\delta_{\rm ar}\).
The shape table is exact rational at the delivered \(c_r\), or is
stored with error at most \(\delta_{\rm ar}\). Thus
\eqref{eq:arithmetic-error} is at most \(e/80\); input perturbation
and preprocessing error are charged once.

\item \textbf{Geometry and exact quadrature.}
Set
\[
 \Delta=\min\left(\frac18,
     \frac{e}{2048(1+B_{\rm loc})(1+\overline V)}\right).
\]
Use guards \(\eta\le\Delta\), base-node errors at most
\(\Delta/(1+\overline V)\), and split that total base-weight allowance
over the retained intervals. Then
\(\varepsilon_b,\varepsilon_u,\varepsilon_w\le\Delta\), so
\eqref{eq:geometry-error} is below \(e/100\).
Choose the outer Gauss order \(n\ge1\) by the strict rational test
\(16\overline V\overline B_E4^{-n}<e/80\).
This is a terminating geometric test, with \(n=O(L^2)\).

\item \textbf{Retained-integral data.}
Set
\[
 \varepsilon_Q=\frac{e}{80[1+(2\overline V+\Delta)K_{\rm loc}]},
\]
and choose real-node upper bounds \(\overline A,\overline R_k\)
for the factors in \eqref{eq:node-data-error}. With
\[
 M_Q=1+2\sum_{k=0}^K
  [\overline R_kD_k+(\overline A+1)(D_k+\overline R_k+1)],
 \quad\delta_{\rm RS}\le\min(1/4,\varepsilon_Q/[4(1+M_Q)]),
\]
deliver each amplitude, power, coefficient, and exact-transfer
approximation to error \(\delta_{\rm RS}\).
Then \eqref{eq:node-data-error} is at most \(\varepsilon_Q/4\)
and \eqref{eq:RS-data-budget} at most \(e/80\).
Use \eqref{eq:contour-delivery}--\eqref{eq:Ck-delivery} for the
coefficients, keeping every \(k\le K\).
The contour order is \(O(L^2+L\log L)\).

\item \textbf{Final circuit.}
Evaluate the finite moment, kernel, and derivative-channel recurrences
at the already delivered inputs. Their input differences have been paid
in the preceding categories. In this final comparison treat those
delivered inputs as exact, and propagate only new primitive and operation
errors using \eqref{eq:roundoff-updates} and the elementary tail bounds.
Refine their dyadic tolerances until the finite positive error sum is
below \(e/160\), reserving another \(e/160\) for final quantization.
Every division has a positive certified separation; the graph is finite,
so its bound tends to zero as the new tolerances tend to zero.
This proves finite termination with \(E_{\rm round}<e/80\).
\end{enumerate}
The exact-rational reference construction has zero internal operation
rounding. For bounded bit work, evaluate individual rational terms,
quantize to a common dyadic grid, and accumulate as integers. A sum of
\(S\) terms costs its actual \(S\) quantization allowances and
\(\lceil\log_2S\rceil\) guard bits. A floating alternative uses
\eqref{eq:roundoff-updates} and \(uS_{\rm lin}\le1/4\).
Provider-internal errors are included in their delivered radii, not
charged again as final-circuit errors.

\begin{theorem}[Effective canonical evaluation]\label{thm:effective}
Fix positive dyadic \(r,b\), and use the canonical source
\eqref{eq:canonical}. For each accepted exact query of
\cref{subsec:input-class} with \(L\ge64\nu\) in
\eqref{eq:linear-onset}, there is a terminating finite algorithm
delivering the field, spatial derivatives through order three, and first
time derivative, with the combined output error
\[
 \frac{\max_{0\le j\le3}|\widehat h_j-\partial_z^jh|
          +|\widehat h_T-\partial_Th|}{\calE_{\mathrm o}(-T,z)}
 \le\e^{-L^2/4}.
\] The same conclusion holds at any smaller height
whose structural and half-target gates are strictly certified.

The construction retains both arms, all \(n\le N\), all
\(C_0,\ldots,C_{2p}\), the exact adjacent-cutoff transfer, and the
bounded real completion integral. It returns a common scale as a
rational mantissa and an unrestricted binary integer exponent.
\end{theorem}
\begin{proof}
\Cref{thm:linear-onset} pays the structural hypotheses and half-target
truncation. \Cref{sec:providers} supplies every delivered quantity in
the preceding finite allocation. Its five evaluation categories are
at most \(e/80<\varepsilon_X/20\), and scale error at most
\(e/8<\varepsilon_X/2\). \Cref{thm:evaluation} now gives the total
target. Each majorant, tail test, and strict descriptor comparison is
effective, and each refinement has been shown to terminate.
\end{proof}

This theorem supplies an analytical algorithm, not a claim that a
complete software implementation has been certified for every input.
Exact coordinates are part of the contract. Uncertain input coordinates
need an additional derivative-based error budget. The exact frozen model
has zero heat residual; separately rounded channels have a defect bounded
by their paid errors. Its spatial derivative requires the mixed \(T,z\)
channel, as explained in \cref{sec:evaluation}.

\subsection{Delivered precision, total bit work, and streaming}
\label{subsec:bit-complexity}
Use elementary multiplication and division of \(B\)-bit integers at
cost \(O(B^2)\). Let \(\mathcal B\) be actual binary input length,
including center, query, and descriptor witnesses, and define
\[
 s_\sigma=\lceil\log_2^+(1/\sigma)\rceil,\qquad
 Q=2+\mathcal B+P+s_\sigma+\lceil L\rceil^2.
\]
Here \(P\ge1\) covers every externally delivered provider tolerance
and final-circuit primitive/quantization request, including large-phase
and full-sum allowances. Supplied descriptor witnesses can be verified
at their actual length; \(s_\sigma\) retains the dependence of finding
them by refinement. The boundary guard construction requires no further
edge-separation parameter.

For fixed chart data and factor-controlled allocations,
\begin{equation}
 P=\frac{3L^2}{8\log2}+O_{r,b}(L\log L)
 \label{eq:effective-precision}
\end{equation}
is sufficient. Indeed all conditioning and positive path costs have
logarithms \(O(L\log L)\), both quadrature orders are \(O(L^2)\),
and \(\log N=O(L)\). This bounds delivered-data and final-circuit
accuracy requests; provider-internal isolation and temporary operand
lengths are counted separately.

\begin{theorem}[Effective work and space]\label{thm:bit-cost}
For the construction in \cref{thm:effective}, an effective
chart-dependent constant \(C(r,b)\) gives the sufficient bounds
\begin{equation}
 T_{\rm bit}\le C(r,b)(N+1)Q^{40},\qquad
 M_{\rm bits}\le C(r,b)Q^{20}.
 \label{eq:total-bit-cost}
\end{equation}
Space uses streaming of arithmetic modes and sequential gamma calls.
General input/tolerance choices retain their actual \(\mathcal B,P\).
\end{theorem}
\begin{proof}
The finite majorants give outer and contour orders \(n,q=O(Q)\)
and \(p,K,J,m=O(Q)\). Gauss isolation depth and physical-node input
length are \(O(Q^4)\); weight numerators and denominators have
\(O(Q^5)\) bits. \Cref{prop:gauss-provider} costs \(O(Q^{16})\)
for the two reusable base rules. Outer-node gamma calls have
\(b_\Gamma=O(Q^4)\), so \cref{prop:gamma-provider} costs
\(O(Q^{24})\) each and \(O(Q^{25})\) over their \(O(Q)\) calls,
with \(O(Q^{12})\) workspace.

Common shape and weighted coefficient preprocessing use
\(O(p^3+K^2)\) scalar operations. Reusable contour samples and their
composition use \(O(n(qK+K^2))\); arithmetic moments and channels use
\(O(Np)\). Every remaining individual input or temporary operand is
bounded by \(O(Q^5)\) bits, including reciprocal powers through
\(O(K)\). The elementary-series construction in \cref{app:elementary}
therefore costs \(O(Q^{30})\) per call. At most
\((N+1)O(Q^5)\) such calls and term operations occur.
Their \(O((N+1)Q^{35})\) work, together with rounding, verification,
and finite error-ledger arithmetic, fits \eqref{eq:total-bit-cost}.

Retain the common polynomial tables, two base Gauss rules, and contour
samples, but process one outer node or arithmetic mode at a time.
The remaining elementary workspace is \(O(Q^{15})\); all tables and
accumulators fit \(O(Q^{20})\). A full linear sum adds only
\(O(\log N)\) integer bits. No product of \(N\) unrelated rational
denominators is accumulated. The scale exponent has \(O(L)\) bits;
its expanded denominator is never materialized. Substitution of
\eqref{eq:explicit-calibration} makes every fixed constant effective.
\end{proof}
The exponents \(40,20\) are deliberately loose upper bounds.
Since \(N\asymp\e^{L/2}\), this is not polynomial total work in
\(L\). Materializing all mode data would also change the storage bound.
No practical runtime at \(X_{\rm new}\), sharp complexity, or arithmetic
compression is asserted. Likewise the compact encoding of \(X_{\rm new}\)
does not make its expanded input length disappear: its binary expansion
has \(128\nu+9\) bits, whereas its exponent has \(O(\log\nu)\) bits.
The earlier \(X_*\) remains a sufficient fallback.

\subsection{Supporting computational evidence}
Finite scalar/gamma and Gauss implementations and certificate replay
support parts of these constructions. The linked computational note
\cite{rec12} contains the exact cases, error witnesses, rejection controls,
resource accounting, preserved failed launch, and reproduction guide;
\cite{rec13} places that evidence in the research sequence.
The analytical theorems above do not depend on those finite tests.
The complete canonical-field evaluator and its software certification
remain unfinished, and separate replay does not supply independent
specialist validation of the mathematical premises.

\section{Further analytical questions}
\label{sec:outstanding}
The fixed-order and growing-order comparisons now lead to an effective
finite construction with explicit calibration and cost. The useful next
advances should strengthen a quantitative conclusion of that theory.
The following order separates analytical goals from the computations
that can support their development.

\begin{enumerate}[leftmargin=*,itemsep=.65em]
\item \textbf{Extend beyond the paid moving-endpoint regime.}
The compact-family theorem fixes \(\tau,B,p\); the detailed growing-order
construction in \cref{thm:growing} fixes \(1\le T\le2\).
The integrated AN08 proof in \cref{sec:moving-endpoint} pays
\(\tau(X)=1/L\), \(B=2\), \(p=\lfloor L/2\rfloor\) at unchanged
accuracy, with sufficient \(L\ge1024\nu^2\), strict query witnesses,
a paid endpoint tag and \(O(L^3)\) sufficient outer nodes
\cite{rec20,rec20review}. Further extensions should declare
\(\tau(X),B(X),p(X)\) and prove simultaneous control of the buffered
factor \(\e^{Y^2/(4\tau)}\), the inverse-time derivative factors,
Gaussian moments, and tails. Beginning with a lower endpoint approaching
zero while keeping the upper endpoint fixed is a narrower problem than
both endpoints moving. The signed crossover, retained-arithmetic positive band and compact
positive atlas in \cite{ansatzcomp} retain their own domains, envelopes
and weaker field allowances. Strong negative precision through zero
remains open. The recurrence comparison supplies no corrected physical
consumer theorem. The source registers own research priority and status.


\item \textbf{Further coefficient and integral improvements.}
\Cref{prop:sharper-rs,prop:smaller-contour} already reduce the uniform
coefficient constants by more than ten and save eleven nodes on each
contour edge at unchanged accuracy. Further work should improve the
order dependence, denominator-safe local delivery, or the actual
retained-integral allowance. The weighted-table factorial bound and
arithmetic-transfer costs remain. A new count or error gain must be
proved before selecting supporting computations; the closed constant
comparison should not be repeated as an open task.

\item \textbf{Refine the remaining onset and parameter choices.}
\Cref{thm:linear-onset} replaces the earlier quartic sufficient
log-height by \(64\nu\) with every original gate intact. A further gain
must compare a proved smaller continuous envelope, or audit simultaneous
changes in window and retained orders. The current linear bound is
sufficient, not optimal; one isolated verified height does not replace
an all-height theorem. Such refinement follows the completed additive-cost
comparison rather than repeating its former coarse-envelope target.

\item \textbf{Quantify any compression or continuous deformation.}
The full \(N\asymp\sqrt X\) arithmetic population dominates the present
cost. A useful compression must reconstruct it with a certified field
and derivative error within the same budget, and must retain the signed
cancellations needed by the application. The local deformation parameter
\(s\) and exact recurrence are starting identities; whole-line convergence,
resummation, and moment closure require separate uniform estimates.
Removing the bounded completion integral likewise needs a paid
reconstruction theorem.

\item \textbf{Transfer the estimates to a specified physical consumer.}
A concrete observable should determine the required norms and channels.
The parent and ansatz companions supply a specialized fixed-P4 compact
switch consumer with its own absolute rates; it does not transfer the
present strong negative-time estimate to an arbitrary observable. A new
consumer must pay regulator dependence, signed covariance, and spatial
and temporal joining defects. State the resulting consumer
bound before selecting supporting computations. An outer-envelope field
estimate and exact frozen heat evolution do not by themselves supply
these quantities.
\end{enumerate}

Independent mathematical validation is a standing need across these
questions. The classical normalization and coefficient identities should
be checked against their primary proofs, and the comparison, calibration,
and cost arguments reviewed independently. Internal LLM assessments and
component replay retain their stated evidence limits.

The remaining software work includes pole-safe coefficient delivery,
scale and large-phase handling, exact frozen geometry, retained-integral
data, and the full mode and derivative circuits. Bounded component
certificates can test those constructions and expose concrete failures.
Completing that implementation is distinct from improving the analytical
onset or proving a new time regime. Its measured performance should not
be used to extrapolate a theorem beyond the implemented cases.

The source-specific scope also remains essential. Euler-factor and
finite-prime controls lack the unchanged canonical coefficients;
weighted-square and cutoff controls lack the complete xi/Riemann--Siegel
initial identity; reflected, smooth, and positive-primitive controls lack
the same complete gamma/arithmetic source \cite{controls,avoid}.
Their heat identities can survive. No RH, simplicity, noncollision,
or physical-transfer conclusion is asserted by the present approximation
and evaluation results.

\appendix
\section{Finite onset conditions and order-dependent constants}
\label{app:onsets}
The following sufficient conditions make the fixed-order theorem
fully parameterized. In addition to the classical real
Riemann--Siegel and sectorial gamma onsets, impose
\begin{subequations}\label{eq:structural-gates}
\begin{align}
 &L\ge16,\quad X\ge\max(1,4(R+Y+1)),\quad V_B\le X/4,\quad
                         X-V_B\ge X/2,\\
 &a(X+V_B)-a(X-V_B)<1,\qquad
                       2p\le (X-V_B)/(4\pi),\\
 &DV_B/X\le1/2,\qquad 2AV_B^2/X\le1,\\
 &D_0/X\le1-\pi/8,\qquad
 B\{(D_0/X)(\pi/4+L/2)+D_0^2/X^2\}\le1,\\
 &V_B\ge\max(2R,W+1,32B).
\end{align}
\end{subequations}
Here \(D_0\) is a fixed constant in the sectorial slope formula
\[
 \lambda_{\epsilon n}=-\pi/8+\epsilon\ii v_n/2+\delta_\epsilon,
 \qquad
 v_n=\log(a_X/n),\qquad \abs{\delta_\epsilon}\le D_0/X.
\]
The optional inherited real-chart restriction
\(V_B\le\log^2X\) is also eventual.
Writing \(\log^+a=\max(0,\log a)\), the fixed-order tail
conditions are
\begin{subequations}\label{eq:fixed-tail-gates}
\begin{align}
 L^2&\ge(P_0+p-1)\log X+\log^+ A_{\mathrm{mod}},\\
 L^2&\ge(p+1)\log X+\log^+ A_{\mathrm{art},p},\\
 X^2/(128B)&\ge\pi X/8+(p-1)\log X+\log^+ A_{\mathrm{far}}.
\end{align}
\end{subequations}
They put each tail in \cref{sec:comparison} below
\(X^{-(p+3/4)}\). Every condition is eventual for fixed
\(p,\tau,B,r,b\). The factors
\(\e^{Y^2/(4\tau)}\), \(\tau^{-1}\), \(\tau^{-2}\),
and the full moments in \cref{eq:gaussian-budgets,eq:tail-moments}
remain in these constants.

\subsection{Common gates at growing order}
\label{app:growing-gates}
For \(1\le p\le cL\), \(B=2,\tau=1\), replace the retained-order
condition in \eqref{eq:structural-gates} by
\(2cL\le(X-V)/(4\pi)\).
Include the fixed outer-calibration gate
\(C_{\mathrm{cal}}(1+L)/X\le1\).
Do not impose \eqref{eq:fixed-tail-gates} at growing order.
Instead, with \(C_{\mathrm o}=C_\chi/c_{\mathrm o}\), the following
four inequalities suffice for \cref{thm:growing}:
\begin{subequations}\label{eq:growing-gates}
\begin{align}
 (p+3/4)\log X
 &\ge(\eta+1/8)L^2+\log^+(4C_{\mathrm o}Q_p),\\
 (16\theta-\eta-1/8)L^2
 &\ge(P_0-7/4)\log X+
                   \log^+(196C_{\mathrm o}A_{\mathrm{mod}}^\theta),\\
 (16\theta-\eta-1/8)L^2
 &\ge\tfrac14\log X+
                   \log^+(196C_{\mathrm o}A_{\mathrm{art},p}^\theta),\\
 X^2/256
 &\ge\pi X/8+(\eta+1/8)L^2+
                   \log^+(196C_{\mathrm o}A_{\mathrm{far}}).
\end{align}
\end{subequations}
Each pays at most one quarter of the target.
For the half-target truncation required in
\cref{sec:evaluation}, replace \(4,196\) by \(8,392\).
These are the operative finite inequalities.
\Cref{thm:explicit-onset} now supplies an explicit sufficient all-height
threshold; a smaller height can be accepted by its finite verifier.
For the moving domain use the distinct constants
\eqref{eq:moving-truncation-constants}, unchanged half-target factors
\(8,392\), and the continuous slack in
\cref{subsec:moving-onset-proof}. Do not substitute the moving endpoint
into the fixed-domain constants above.

\section{Evolved population bounds and finite level size}
\label{app:population}
The error proof used the real initial population. One can also
bound the entire evolved population without discarding the
near-cutoff modes. Put \(a_0=\pi/8\), \(c_Y=(Y-1)/2\),
\(c_Y^+=\max(c_Y,0)\), and for an integer \(d_0\ge0\) define
\begin{subequations}\label{eq:Pi}
\begin{align}
 W_{d_0}(\tau)
 &=2^{d_0}+\frac{3^{d_0}}2\sum_{k=0}^{d_0}\binom{d_0}{k}
              (8/\tau)^{(k+1)/2}\Gamma((k+1)/2),\\
 \Pi_{d_0}(\tau,B)
 &=4(4\pi)^{-1/4}4^{d_0}
   \exp\{R+Y+Ba_0^2+1+(Y+1)/2+2(c_Y^+)^2/\tau\}
                   W_{d_0}(\tau).
\end{align}
\end{subequations}
The slope gates in \eqref{eq:structural-gates} imply
\[
 \abs{\e^{T\lambda_{\epsilon n}^2}}
       \le\e^{Ba_0^2+1}\e^{-\tau v_n^2/4},
 \qquad
 \abs{\e^{\lambda_{\epsilon n}d}}
       \le\e^{R+Y}\e^{Yv_n/2}.
\]
Also \(1+\abs\lambda\le4(1+v_n)\).
The shell \(j\le v_n<j+1\) contains at most \(2a_X\e^{-j}\)
integers. This remains true in the last nonempty shell,
where \(a_X\e^{-j}\ge1\).
In that shell \(n^{-1/2}\le a_X^{-1/2}\e^{(j+1)/2}\).
Use
\[
 c_Yj-\tau j^2/4
       \le2(c_Y^+)^2/\tau-\tau j^2/8.
\]
For \(j\ge1\), compare with the integral on \([j-1,j]\);
\(j+2\le3(v+1)\) there and
\(\e^{-\tau j^2/8}\le\e^{-\tau v^2/8}\).
The half-Gaussian moments give exactly the bound \(W_{d_0}\).
Summing both arms proves
\begin{equation}
 \sum_{\epsilon,n\le N}
 \abs{b_{\epsilon n}\e^{\lambda_{\epsilon n}d+T\lambda_{\epsilon n}^2}}
       (1+\abs{\lambda_{\epsilon n}})^{d_0}
       \le \Pi_{d_0}(\tau,B)gX^{1/4}.
 \label{eq:evolved-population}
\end{equation}

For the moment polynomial define
\[
 Z_m(B)=m!\sum_{\nu=0}^{\lfloor m/2\rfloor}
       \frac{B^\nu(W+2B)^{m-2\nu}}{\nu!(m-2\nu)!}.
\]
Since \(\abs{d+2T\lambda}\le(W+2B)(1+\abs\lambda)\),
the whole \(m\)-th moment sum is at most
\(Z_m\Pi_m gX^{1/4}\).
For \(j\ge1\) this gives
\begin{equation}
 \sum_{\epsilon=\pm1}\abs{U_{\epsilon,j}}
 \le gX^{1/4-j}
       \sum_{k=1}^j\frac{A^kD^{j-k}}{k!}\binom{j-1}{k-1}
                                   Z_{j+k}\Pi_{j+k}.
 \label{eq:level-size}
\end{equation}
For \(\tau=1/L\), the degree-dependent comparison
\eqref{eq:moving-evolved-ratio} is the separate evolved-population
ledger; the order-independent truncation multiplier
\eqref{eq:moving-truncation-ratio} does not bound it.
On the slab \([1,2]\), the explicit factorial sums give fixed
\(K_0,K_1\) such that \(Z_m\Pi_m\le K_0K_1^m(m!)^2\).
For example, \(Z_m/m!\) is a coefficient of
\(\exp((W+4)s+2s^2)\), so \(Z_m\le m!\e^{W+6}\);
the half-gamma bound similarly controls \(\Pi_m\).
Consequently the \(j\)-th level has the absolute envelope
\[
 gX^{1/4-j}\exp(O((j+1)\log(j+2))).
\]
This is a bound for the entire population, not a signed leading
asymptotic, a monotonicity statement, or a convergence theorem
for the sum over all levels.

\section{Sources, scope, and manuscript correspondence}
\label{app:provenance}
Internal note paths in the bibliography are relative to
\nolinkurl{papers/investigations-5/01_heat_density_asymptotics/}.
The analytical inputs and their derived uses are deliberately
separated:
\begin{center}
\small
\begin{tabular}{@{}>{\raggedright\arraybackslash}p{.26\textwidth}>{\raggedright\arraybackslash}p{.34\textwidth}>{\raggedright\arraybackslash}p{.29\textwidth}@{}}
\toprule
Input or research note & Use here & Scope retained\\
\midrule
Polymath Proposition 6.2 &
Real canonical remainder for each positive \(K\) &
No evolved negative-time theorem imported\\[3pt]
DLMF gamma asymptotics &
Sectorial derivative and calibration bounds &
Fixed sectors, explicit eventual gates\\[3pt]
Arias exact identities &
Weighted coefficients and \(F_0\) derivatives &
Preliminary precision conclusions excluded\\[3pt]
DLMF Gauss properties &
Positivity and degree \(2n-1\) exactness &
Complex error bound derived here\\[3pt]
Recurrence Notes 02--04 &
Finite hierarchy and compact comparison &
Fixed \(p,\tau,B\) first\\[3pt]
Recurrence Note 05 &
Simultaneous order costs and outer accuracy &
\(p=O(L)\) on \([1,2]\)\\[3pt]
Recurrence Note 06 &
Regularity, conditioning, and evaluation budget &
Modular provider and arithmetic contracts\\[3pt]
Recurrence Note 09 &
Finite elementary, scale, coefficient and guard constructions &
Strict descriptors and exact queries\\[3pt]
Recurrence Note 10 &
Shifted gamma, Gauss root counts, total bit work and storage &
All modes and actual separation retained\\[3pt]
Recurrence Note 11 &
Uniform constants and explicit all-height onset &
Fixed radii, \(1\le T\le2\)\\[3pt]
Recurrence Note 12 &
Linked finite component replay &
Supporting evidence only\\[3pt]
Recurrence Note 16 &
Additive-cost linear sufficient log-height &
All original gates and half-target budget\\[3pt]
Recurrence Note 18 &
Sharper uniform RS bounds and contour count &
Same domains and delivered accuracy\\[3pt]
Recurrence Note 20 &
Integrated moving-endpoint proof in \cref{sec:moving-endpoint} &
\(L^{-1}\le T\le2\), strict queries and paid tag\\[3pt]
Book manuscript &
Common notation and cross-region synthesis &
Proof ownership remains here\\[3pt]
Ansatz Notes 17--20 and companion &
Positive finite hierarchy and compact bridge &
Fixed positive order and its own paid domains\\
\bottomrule
\end{tabular}
\end{center}
The preceding ansatz results in \cref{subsec:retention} are
retained as separately sourced theorems with their original
domains. Their first-correction formula is reconciled with
the present one in \cref{subsec:low-orders}; no differentiation
rule is transferred from one construction to the other.
For the current positive comparison use
\cref{subsec:positive-comparison} and the proof labels in
\cite{ansatzcomp}; the Book's negative-time contracts point to
\cref{thm:fixed,thm:growing,thm:linear-onset,thm:effective,thm:bit-cost},
\cref{thm:moving-endpoint-delivery} and
\cref{prop:sharper-rs,prop:smaller-contour} here.
Note 19 retains AN08's original request. The later Note 20/Review 14
return is Complete/CLOSED in the live register; its full proof is now
integrated in \cref{sec:moving-endpoint}
\cite{rec19,rec20,rec20review,recstatus}. The input interface and
provider cost changes are given in \cref{subsec:moving-input,subsec:moving-cost}.
The repository result index and control screen
\cite{results,controls}, together with its scoped-error cautions
\cite{avoid}, govern the interpretation.

The source notes were developed and internally assessed with
gpt-6.1-sol, recorded effort ultra. This manuscript was drafted
for Edward Baker with gpt-6-astra, recorded effort max;
the linear-onset and uniform RS-bound strengthenings use gpt-6.1-sol,
actual effort ultra. The Book-companion reconciliation and subsequent moving-endpoint proof
integration were prepared with gpt-6.1-sol, actual effort ultra.
The mathematical status remains that of the stated proofs
from named inputs. The modular contracts are supplied by finite
constructions on the declared input class; complete software
certification and independent specialist review remain open.
Compilation, reference checks, and visual inspection concern
document production. They are not numerical evidence for
the analytical estimates.

\section{Elementary enclosures and alternative gamma construction}
\label{app:elementary}
This appendix records the finite tails behind \cref{lem:elementary}.
For positive rational \(x=2^ky\), \(1\le y<2\), put
\(t=(y-1)/(y+1)\), so \(\abs t\le1/3\). Then
\begin{equation}
 \log y=2\sum_{j=0}^{M-1}\frac{t^{2j+1}}{2j+1}+R_M,
 \qquad\abs{R_M}\le
 \frac{2\abs t^{2M+1}}{(2M+1)(1-t^2)}.
 \label{eq:log-series-provider}
\end{equation}
Use \(t=1/3\) for \(\log2\), with its error multiplied by \(\abs k\).
Positive interval logarithms use monotonicity and separated endpoints.
Square roots use exact square comparisons and bisection after scaling
to \([1,4]\).

For \(\abs t\le1/2\), the first \(M\) arctangent terms have
remainder at most \(\abs t^{2M+1}/(2M+1)\). Reduce
\(1/2\le t\le1\) by
\(\arctan t=\arctan(1/2)+\arctan((t-1/2)/(1+t/2))\),
whose new argument is at most \(1/3\); use oddness and
\(\arctan t=\pi/2-\arctan(1/t)\) for the other ranges.
Tangent addition and the angle range give
\(\pi=4[\arctan(1/2)+\arctan(1/3)]\).
For \(x>0\),
\(\Log(x+\ii y)=\log(x^2+y^2)/2+\ii\arctan(y/x)\).
Rational reduction boundaries can be subdivided; either valid neighboring
formula agrees at the boundary.

Enclose \(a/\log2\) to width at most \(1/8\) and round its rational
midpoint to an integer \(k\). This is an encoding choice. Refinement
with that \(k\) fixed certifies \(\abs{a-k\log2}<2\).
For a reduced real argument of modulus at most two, the degree-\(M\)
Taylor error is at most \(9\,2^{M+1}/(M+1)!\).
For phase reduction choose a nearby multiple of \(2\pi\) from an
enclosure of width \(1/8\); its residual has modulus at most
\(9\pi/8<4\). The corresponding complex Taylor error is at most
\(81\,4^{M+1}/(M+1)!\).
Refine \(\pi\) and \(\log2\) with the actual chosen multipliers.
Relative real-exponential error \(\tau\) is obtained by absolute
Taylor error at most \(\tau/9\), since the reduced exponential
is at least \(1/9\), together with a separate input-error allowance.
For complex input perturbation use
\[
 \abs{\e^{z+q}-\e^z}\le\e^{\Rea z}(\e^{\abs q}-1).
\]
If \(\abs a\ge d_0>0\) and \(\abs{\widehat a-a}\le\eta<d_0\),
then the reciprocal error is at most \(\eta/[d_0(d_0-\eta)]\).

Every truncation is chosen by a strict rational tail test. For a reduced
series with \(b\)-bit rational input and requested \(P\)-bit accuracy,
\(O(P+b)\) terms suffice. Explicit common denominators such as
\(q^{2M+1}(2M+1)!\) have \(O((P+b+10)^2)\) bits.
Elementary integer arithmetic gives the conservative
\(O((P+b+10)^6)\) work bound used in \cref{thm:bit-cost}.
This refers to reduced arguments and mantissa accuracy. Absolute
accuracy after rescaling an exponential additionally counts its output
magnitude; there is no polynomial absolute-error assertion for an
unrestricted enormous exponential.

\subsection{Euler-product fallback and its internal error}
\label{app:euler-provider}
The elementary construction of \cite{rec09} remains useful as a second
finite formula. Let \(\zeta=(1+\ii u)/4\), \(v=\zeta-1\),
choose rational \(B_\Gamma\ge\abs v\), and take integer
\(s\ge\max(2B_\Gamma,2)\). Set \(q=B_\Gamma/(s+1)\le1/2\) and
\[
 F_s(v)=v\log s-\sum_{j=1}^s\Log(1+v/j).
\]
Euler's product \cite{dlmfgammaproduct} and the harmonic-sum remainder
give
\begin{align}
 \abs{F_s-\log\Gamma(1+v)}
   &\le B_\Gamma/s+B_\Gamma^2/[2s(1-q)],\nonumber\\
 \abs{F_s'-(\log\Gamma)'(1+v)}
   &\le1/s+B_\Gamma/[s(1-q)],\nonumber\\
 \left|\frac{F_s^{(m)}-(\log\Gamma)^{(m)}(1+v)}{m!}\right|
   &\le\frac1{m(m-1)(1-q)^m s^{m-1}}\quad(m\ge2).
 \label{eq:euler-fallback-tails}
\end{align}
To verify these inequalities subtract the locally convergent logarithmic
product. Its linear coefficient differs by
\(H_s-\log s-\gamma\in[0,1/s]\); the remaining log terms obey
\(\abs{\Log(1+v/j)-v/j}\le\abs{v/j}^2/[2(1-q)]\).
Differentiate the tail and use
\(\sum_{j>s}j^{-m}\le s^{1-m}/(m-1)\).
The finite derivative formulas are
\[
 F_s'=\log s-\sum_{j=1}^s(j+v)^{-1},\qquad
 F_s^{(m)}/m!=\frac{(-1)^m}{m}\sum_{j=1}^s(j+v)^{-m}.
\]
All denominators have real part at least \(1/4\).
A common cutoff
\(s\ge\max(2B_\Gamma,2,2(1+B_\Gamma+B_\Gamma^2)/\delta)\)
makes every truncation error at most \(\delta\).

For desired output error \(\eta\), choose \(\delta=\eta/4\).
Deliver each value logarithm to
\(\eta/[8(s+1)(B_\Gamma+1)]\), \(\log s\) to
\(\eta/[8(B_\Gamma+1)]\), each normalized reciprocal-power summand
to \(\eta/[8(s+1)]\), and the derivative's \(\log s\) to
\(\eta/8\). Exact rational summation then leaves output slack.
This takes \(O((1+B_\Gamma^2)/\delta)\) terms. Using only
output-level unit roundoff in this much longer sum need not meet the
accumulation guard. Termwise precision or exact summation is essential;
\cref{prop:gamma-provider} subsequently improves the term count.

\section{Explicit conditioning for Gauss data}
\label{app:gauss-conditioning}
For \(s\ge1\) let \(A_s=2^ss!P_s\in\mathbb Z[x]\), with leading
coefficient \(a_s=(2s)!/s!\), and let \(c_s\) be that of \(P_s\).
Define
\begin{align*}
 \sigma_s&=[a_s^{2s-1}2^{s(s-1)}]^{-1},&
 d_s&=c_s\sigma_s^{s-1},\\
 S_1&=\sum_jj\abs{[x^j]P_s},&
 S_2&=\sum_jj(j-1)\abs{[x^j]P_s},\\
 J_*&=\sum_j\abs{[x^j]J_s},&
 J_1&=\sum_jj\abs{[x^j]J_s},\\
 L_w&=2J_1/d_s+4J_*S_2/d_s^2,&
 \mu_s&=[4\max(1,2sS_1/d_s)]^{-1}.
\end{align*}
For \(s\ge2\), the nonzero integer resultant of \(A_s,A_s'\)
equals, up to sign, \(a_s^{2s-1}\) times the squared-gap product.
All gaps are at most two, so every gap is at least \(\sigma_s\).
For \(s=1\) there is no gap condition. Likewise the nonzero integers
\(A_s(\pm1)\) give endpoint distances at least
\(1/(a_s2^{s-1})\). At a root \(\abs{P_s'}\ge d_s\).
On an interval of width at most \(d_s/[2\max(1,S_2)]\), its
derivative has modulus at least \(d_s/2\); differentiation of
\(J_s/P_s'\) gives Lipschitz bound \(L_w\).

For the nodal Lagrange polynomial,
\(\abs{l_i'}\le sS_1/d_s\) and \(l_i(\xi_i)=1\).
Integrating its square toward the farther endpoint gives
\(\omega_i\ge\mu_s\). Thus it suffices to isolate to width
\begin{equation}
 \delta_*\le\min\left\{\frac18,\frac{\sigma_s}4,
 \frac{d_s}{4\max(1,S_2)},\eta_x,
 \frac{\eta_w}{2s\max(1,L_w)},
 \frac{\mu_s}{2\max(1,L_w)}\right\}.
 \label{eq:gauss-isolation-target}
\end{equation}
Take \(2^{1-h}\le\delta_*\). All conditioning lengths are
\(O(s^3\log(s+2))\), giving the depth in \cref{prop:gauss-provider}.

For weight arithmetic form
\[
 Q_s(c)=s!\int_{-1}^1\frac{A_s(t)-A_s(c)}{t-c}\dd t,
 \qquad
 [c^r]Q_s=2\sum_{\substack{j>r\\j-r\ {
m odd}}}
                   [t^j]A_s\frac{s!}{j-r}.
\]
Both integer polynomials have \(O(s\log(s+2))\)-bit coefficients.
The rational weight at a midpoint is
\(Q_s(\widehat\xi)/(s!A_s'(\widehat\xi))\).
At \(\widehat\xi=c/q\), homogeneous evaluation cancels the common
\(q^{s-1}\), giving \(O(sU_s)\)-bit numerators and denominators.
This proves the weight-size assertion used in the total cost ledger.

The global depth is sufficient, not mandatory. An isolating interval
with a certified nonzero \(P_s'\) can instead use direct interval
evaluation of \(J_s/P_s'\). If its node and true-weight enclosures
already meet the assigned radii and positivity tests, it is a valid
earlier stopping certificate for the same algorithm. This observation
does not improve the stated worst-case complexity; its bounded software
use is documented in the supporting note \cite{rec12}.

\begin{thebibliography}{99}
\small\raggedright
\bibitem{polymath}
D. H. J. Polymath,
\emph{Effective approximation of heat flow evolution of the Riemann
\(\xi\) function, and a new upper bound for the de Bruijn--Newman constant},
arXiv:1904.12438v2 (4 August 2019), especially Proposition 6.2.
\url{https://arxiv.org/abs/1904.12438v2}.

\bibitem{dlmfgamma}
NIST Digital Library of Mathematical Functions,
\emph{Section 5.11: Gamma Function Asymptotic Expansions},
especially (5.11.1)--(5.11.3) and the sectorial remainder discussion.
\url{https://dlmf.nist.gov/5.11}.

\bibitem{dlmfgauss}
NIST Digital Library of Mathematical Functions,
\emph{Section 3.5(v): Gauss Quadrature},
especially (3.5.18)--(3.5.20\_1).
\url{https://dlmf.nist.gov/3.5.v}.

\bibitem{ariasI}
J. Arias de Reyna,
\emph{High precision computation of Riemann's zeta function by the
Riemann--Siegel formula, I},
Mathematics of Computation \textbf{80} (2011), 995--1009.
\url{https://doi.org/10.1090/S0025-5718-2010-02426-3}.
Primary proof not independently rechecked in this project;
the real remainder is imported through Polymath.

\bibitem{ariasII}
J. Arias de Reyna,
\emph{High Precision Computation of Riemann's Zeta Function by the
Riemann--Siegel Formula, II}, arXiv:2201.00342v1 (2 January 2022).
\url{https://arxiv.org/abs/2201.00342v1}.
Selected exact identities (29), (62), (86)--(89);
the preliminary implementation/precision conclusions are not used.

\bibitem{ans03}
E. Baker, with LLM assistance,
\emph{Outer correction and a bounded signed-time real-axis bridge},
ansatz research Note 03, 25 October 2026.
Repository: \nolinkurl{ansatz/notes/03_OUTER_CORRECTION_AND_SIGNED_REAL_BRIDGE_20261025.md}.

\bibitem{ans04}
E. Baker, with LLM assistance,
\emph{Fixed-cutoff complex remainder and buffered signed-time crossover},
ansatz research Note 04, 25 October 2026.
Repository: \nolinkurl{ansatz/notes/04_COMPLEX_STRIP_REMAINDER_AND_BUFFERED_CROSSOVER_20261025.md}.

\bibitem{ans05}
E. Baker, with LLM assistance,
\emph{Fixed negative times: a short natural-scale model and an outer-scale representation},
ansatz research Note 05, 25 October 2026.
Repository: \nolinkurl{ansatz/notes/05_FIXED_NEGATIVE_TIME_REPRESENTATIONS_20261025.md}.

\bibitem{ans10}
E. Baker, with LLM assistance,
\emph{First negative-time evolution correction with a retained RS integral},
ansatz research Note 10, 25 October 2026.
Repository: \nolinkurl{ansatz/notes/10_FIRST_NEGATIVE_TIME_EVOLUTION_CORRECTION_20261025.md}.

\bibitem{ans12}
E. Baker, with LLM assistance,
\emph{A fixed compact-negative-time family with endpoint-dependent allowances},
ansatz research Note 12, 25 October 2026.
Repository: \nolinkurl{ansatz/notes/12_COMPACT_NEGATIVE_TIME_FAMILY_20261025.md}.

\bibitem{rec02}
E. Baker, with LLM assistance,
\emph{A paid second-level fluctuation comparison},
recurrence research Note 02, 26 October 2026.
Repository: \nolinkurl{recurrence/notes/02_FIRST_FINITE_FLUCTUATION_COMPARISON_20261026.md}.

\bibitem{rec03}
E. Baker, with LLM assistance,
\emph{Every fixed fluctuation order with matched RS completion},
recurrence research Note 03, 26 October 2026.
Repository: \nolinkurl{recurrence/notes/03_FIXED_ORDER_FLUCTUATION_HIERARCHY_20261026.md}.

\bibitem{rec04}
E. Baker, with LLM assistance,
\emph{Fixed compact-time recurrence family, every fixed order},
recurrence research Note 04, 26 October 2026.
Repository: \nolinkurl{recurrence/notes/04_COMPACT_TIME_RECURRENCE_FAMILY_20261026.md}.

\bibitem{rec05}
E. Baker, with LLM assistance,
\emph{Controlled order growth, retained tails, and outer accuracy},
recurrence research Note 05, 26 October 2026.
Repository: \nolinkurl{recurrence/notes/05_CONTROLLED_ORDER_GROWTH_AUDIT_20261026.md}.

\bibitem{rec06}
E. Baker, with LLM assistance,
\emph{Coefficient, rounding and retained-RS quadrature errors},
recurrence research Note 06, 26 October 2026.
Repository: \nolinkurl{recurrence/notes/06_COEFFICIENT_ROUNDING_AND_RS_QUADRATURE_20261026.md}.

\bibitem{rec09}
E. Baker, with LLM assistance,
\emph{Finite certified providers for the frozen hierarchy},
recurrence Note 09, 26 October 2026.
Repository: \nolinkurl{recurrence/notes/09_FINITE_PROVIDER_ALGORITHMS_20261026.md}.

\bibitem{rec10}
E. Baker, with LLM assistance,
\emph{Polynomial provider bounds, total work, and onset dependencies},
recurrence Note 10, 26 October 2026.
Repository: \nolinkurl{recurrence/notes/10_EFFICIENT_PROVIDERS_AND_TOTAL_COST_20261026.md}.

\bibitem{rec11}
E. Baker, with LLM assistance,
\emph{Explicit uniform calibration and an all-height onset},
recurrence Note 11, 26 October 2026.
Repository: \nolinkurl{recurrence/notes/11_EXPLICIT_CALIBRATION_AND_HEIGHT_ONSET_20261026.md}.

\bibitem{rec12}
E. Baker, with LLM assistance,
\emph{Provider implementation and replay}, recurrence Note 12,
26 October 2026.
\href{run:../../../papers/investigations-5/01_heat_density_asymptotics/recurrence/notes/12_PROVIDER_IMPLEMENTATION_REPLAY_20261026.md}{Linked supporting note: finite cases, certificates, limitations, and reproduction.}
The link resolves from the repository's reading PDF in
\nolinkurl{output/pdf/heat-density-recurrence/}.

\bibitem{rec13}
E. Baker, with LLM assistance,
\emph{Results of recurrence chat0003: providers, costs, calibration and finite replay},
recurrence Note 13, 26 October 2026.
\href{run:../../../papers/investigations-5/01_heat_density_asymptotics/recurrence/notes/13_CHAT0003_RESULTS_SUMMARY_20261026.md}{Linked results summary and evidence boundaries.}

\bibitem{rec16}
E. Baker, with LLM assistance,
\emph{A linear sufficient all-height onset}, recurrence Note 16,
26 October 2026.
\href{run:../../../papers/investigations-5/01_heat_density_asymptotics/recurrence/notes/16_SHARPER_ALL_HEIGHT_ONSET_20261026.md}{Linked proof: additive constants, every gate, and quantitative comparison.}

\bibitem{rec18}
E. Baker, with LLM assistance,
\emph{Uniform RS bounds and fewer contour nodes}, recurrence Note 18,
26 October 2026.
\href{run:../../../papers/investigations-5/01_heat_density_asymptotics/recurrence/notes/18_UNIFORM_RS_BOUND_STRENGTHENING_20261026.md}{Linked proof: full rectangles, delivery errors, and uniform node-count comparison.}

\bibitem{dlmfbernoulli}
NIST Digital Library of Mathematical Functions,
\emph{Bernoulli polynomial Fourier expansion}, (24.8.1).
\url{https://dlmf.nist.gov/24.8.E1}.

\bibitem{dlmflegendre}
NIST Digital Library of Mathematical Functions,
\emph{Legendre recurrences and derivative identity},
Table 18.9.1 and (14.10.5).
\url{https://dlmf.nist.gov/18.9.T1};
\url{https://dlmf.nist.gov/14.10.E5}.

\bibitem{dlmfbeta}
NIST Digital Library of Mathematical Functions,
\emph{Euler's beta integral}, (5.12.1).
\url{https://dlmf.nist.gov/5.12.E1}.

\bibitem{dlmfhalf}
NIST Digital Library of Mathematical Functions,
\emph{Half-gamma modulus}, (5.4.4).
\url{https://dlmf.nist.gov/5.4.E4}.

\bibitem{dlmfzeta}
NIST Digital Library of Mathematical Functions,
\emph{Finite-sum and fractional-part integral for zeta}, (25.2.8).
\url{https://dlmf.nist.gov/25.2.E8}.

\bibitem{dlmfgammaproduct}
NIST Digital Library of Mathematical Functions,
\emph{Gamma infinite products}, (5.8.1)--(5.8.2).
\url{https://dlmf.nist.gov/5.8}.

\bibitem{results}
Shifted zeta positivity project,
\emph{Result summaries}, particularly the heat-flow cluster.
Repository: \nolinkurl{result-summaries/README.md} and
\nolinkurl{result-summaries/11_newman_heat_flow_reductions_and_certificates.md}.

\bibitem{controls}
Shifted zeta positivity project,
\emph{Collision controls and screening}.
Repository:
\nolinkurl{result-summaries/12_collision_controls_and_screening.md}.

\bibitem{avoid}
Shifted zeta positivity project,
\emph{Patterns to avoid}, particularly patterns 01, 02, 05, 07,
12, and 16: target reformulations, shared premises, review status,
scoped exclusions, separate error bounds, and unchecked imports.
Repository: \nolinkurl{avoid/README.md}.
\end{thebibliography}

\end{document}
